\documentclass[../main.tex]{subfiles}
\begin{document}
%==========================================TIÊU ĐỀ TẬP========================================
\begin{table}[h]
    \begin{adjustwidth}{-1cm}{}
        \begin{tabular}{>{\centering\arraybackslash}p{6cm}|>{\raggedright\arraybackslash}p{14cm}}
            \multirow{5}{6cm}
            % Cột bên trái: ÉPISODE và số tập
            {\\[-40pt]
                \flaregothic\fontsize{20pt}{20pt}\selectfont \centering
                \color{eptype}HISTOIRE\color{black}\\
                \gangofthree\fontsize{72pt}{72pt}\selectfont
                \color{epnum}7
                \\[18pt]
                % \flaregothic\fontsize{14pt}{14pt}\selectfont
                % \color{epattr}BIENVENUE, \uline{\color{Banpire} Mel} !
                % \flaregothic\fontsize{18pt}{18pt}\selectfont
                % \color{epattr}ÉPISODE PERDU
                \color{black}
            }
             &
            % Dòng thứ nhất: tên tập tiếng Nhật
            {\fotskip\fontsize{18pt}{18pt}\selectfont \ruby{新}{しん}\ruby{年}{ねん}\ruby{早}{はや}々のテスト！？}
            \\[6pt]
            % Dòng thứ hai: tên tập tiếng Anh
             & {\cafeteria\fontsize{18pt}{18pt}\selectfont A Math Test on the Third Day!?}             \\[4pt]
            % Dòng thứ ba: tên tập tiếng Việt
             & {\vnmsans\fontsize{18pt}{18pt}\selectfont Mùng Ba Tết Thầy - Bài kiểm tra đầu năm!?}    \\[8pt]
            % Dòng thứ tư: Nhân vật xuất hiện
             & {\brian{\fontsize{14pt}{14pt}\selectfont Track used: Lunar Zoo Year - Ultimate Battle}}
        \end{tabular}
    \end{adjustwidth}
\end{table}
%===========================================NỘI DUNG==========================================
\color{black}

\begin{center}
    \uline{8:00 sáng\\Trường Đại học Xây dựng Kawachi.}
\end{center}

Hôm nay là mùng Ba Tết ở Kawachi, và tiết trời đầu xuân dường như đang trong giai đoạn
đẹp nhất của mình. Những cơn gió se lạnh nhẹ nhàng vuốt ve từng tán cây ven đường,
còn nắng mai thì nhuộm vàng sân trường rộng lớn.

Đối với những người sống nơi đây, câu nói ``Mùng Ba tết Thầy'' đã trở thành một nét văn
hóa sâu sắc. Ngày này không chỉ dành để tri ân những người thầy cô giáo đã dìu dắt mình
mà còn là dịp để tôn vinh tri thức, thứ mà người Etsunan (nhất là đất Bắc) xưa nay luôn
coi trọng như báu vật.

Thế nhưng, vào buổi sáng đặc biệt này, sân trường vẫn vắng lặng như tờ. Các bạn học của
Kuon chưa ai quay trở lại sau kỳ nghỉ, còn các thầy cô thì phần lớn vẫn đang tận hưởng
những giây phút sum vầy bên gia đình. Với một người như Kuon, việc không có ai xung
quanh không đồng nghĩa với sự nhàn rỗi. Thay vì dành thời gian để nghỉ ngơi hay rong
chơi như nhiều người khác, cậu chọn cách khai bút đầu xuân bằng một hoạt động độc đáo
nhưng cũng không kém phần\dots\ ``khắc nghiệt.''

Vâng, chính là một bài kiểm tra Toán.

Hikawa Kuon, trong vai trò ``giám thị bất đắc dĩ'', đã tự tay biên soạn một đề thi khó nhằn
để kiểm tra kiến thức cơ bản của mọi người. Không chỉ dừng lại ở các cô gái \hll,
cậu còn gọi cả em gái mình tham gia, tạo thành một ``hội đồng kiểm tra đầu xuân'' đầy
bất ngờ.

Ở vùng đất Bắc của Etsunan, người ta từ lâu đã xem việc học là gốc rễ của mọi
thành công. Những kỳ thi nghiêm ngặt, dù là vào dịp đầu năm, cũng được xem như cơ
hội để thử thách bản thân và đặt ra mục tiêu mới. Và tất nhiên, Kuon của chúng ta không
thể đứng ngoài truyền thống ấy.

Sự việc bắt đầu từ buổi sáng ngày hôm đó\dots

\begin{center}
    \uline{Hồi tưởng - Một ngày trước\\Tầng ba - Nhà Hikawa.}
\end{center}

Tại tầng ba của ngôi nhà, nơi mọi người đang tập trung để bàn kế hoạch cho ngày mùng
Ba Tết. Đây cũng là ngày cuối cùng các cô gái được trải nghiệm không khí Tết ở thủ đô
Kawachi, nên tất cả đều rất háo hức.

Kuon đứng ở trung tâm, với gương mặt điềm tĩnh nhưng ánh mắt sáng rực. Cậu nhẹ nhàng
hắng giọng để thu hút sự chú ý: ``Được rồi mọi người, anh có một vài ý tưởng cho ngày
thứ tư trải nghiệm lễ hội xuân này.''

``Chúng ta sẽ đi dự lễ hội à!?'' Matsuri hớn hở hỏi, đôi mắt của cô sáng rực lên. Cậu con
trai nhíu mày: ``Ờm\dots\ gần đúng. Nhưng hoạt động đó sẽ để dành cho buổi chiều. Buổi
sáng chúng ta có một chút đặc biệt hơn.''

``Vậy thì nó là gì vậy?'' Ayame hỏi, không giấu được sự tò mò.

Kuon hắng giọng một tiếng, trước khi tiếp tục vạch ra kế hoạch của mình, bằng cách dẫn
dắt: ``Trước hết, mấy đứa có biết ý nghĩa của mùng Ba là gì không?''

Mel cố nhớ lại những gì mà cô tìm hiểu về phong tục ngày Tết từ những ngày trước khi
sang đây: ``Theo em đọc trên mạng thì nó là ngày chúc Tết các thầy cô giáo\dots''

A-chan gật gù: ``Đúng rồi đấy, Yozora-san. Ngày Mùng Ba là dịp để tri ân những người
đã dạy mình kiến thức, không chỉ trong trường lớp mà còn cả trong cuộc sống.''

Cậu con trai tiếp lời cô: ``Ở xứ Etsunan này, người ta hay ví thầy cô như `người lái đò',
đưa học trò qua sông tri thức để đến bến bờ tương lai.''

Tất cả các cô gái đều có vẻ như hiểu ý nghĩa về ngày đặc biệt này, theo lời giải thích của
hai người quản lý. Fubuki tươi tỉnh nhìn họ và hỏi: ``Vậy nếu nói như thế, chẳng phải
A-chan và Kuon-senpai cũng được xem là thầy cô sao?''

Nghe vậy, cả A-chan và Kuon đều khẽ khựng lại, tỏ vẻ hơi ngạc nhiên từ lời cô nói, nhất
là từ người được coi là ``quậy banh nóc nhà'' nhất của nhà Tam Giác Xanh.

\begin{dialogue}
    \speak{Kuon} Anh thấy khá ngạc nhiên là một đứa lắm chiêu như em lại nói một câu như thế đấy\dots
    \speak{A-chan} Mấy đứa\dots\ cũng có thể hiểu như vậy, nhưng mà không cần phải khách sáo
    đâu.
    \speak{Kuon} Mấy đứa không cần phải nói quá đâu. Với lại, anh nghĩ thay vì cảm ơn anh
    hay chị A-san, mấy đứa nên tri ân chú CEO của chúng ta trước thì hợp lý
    hơn.
\end{dialogue}

Sora lên tiếng đồng tình với họ: ``Đúng rồi đó. Chính ông ấy là người đã giúp chúng ta
trở thành idol, biến ước mơ của chúng ta thành hiện thực đó.''

Matsuri tỏ vẻ bất ngờ với lời bình của cô ấy: ``Ể? Nhưng nếu chiếu theo định nghĩa mà
Kuon-senpai và A-chan nói, chẳng phải những người dạy học ở trường lớp mới được gọi
là thầy cô thôi chứ?''

Cậu Tổng mỉm cười, tiếp tục giải thích: ``Về bản chất, bất cứ ai dạy mình điều gì có giá
trị thì đều có thể xem là thầy cô cả. Định nghĩa này áp dụng không chỉ cho trường học
mà còn cho gia đình, bạn bè, và cả những người đã giúp mình trưởng thành.''

Kuon ngừng lại, rồi bất ngờ quay sang nhìn Choco với nụ cười ranh mãnh nhưng cũng
hiền từ: ``Đặc biệt là em đó, Choco-sensei.''

Nàng Y tá tỏ vẻ bối rối khi bất ngờ được cậu chỉ mặt gọi tên: ``Hả? Em á? Nhưng em chỉ
là giáo viên y tá thôi mà.''

``Anh để ý em cũng là người thực hiện khảo sát nhiều bài kiểm tra đánh giá năng lực tư
duy cho các thành viên lắm đó. Nên em về cơ bản cũng là giáo viên đấy thôi!'' cậu đáp
lại, rồi hướng về mọi người, phấn chấn nói: ``Mấy đứa, chúc cô giáo của mình đi nào!''

Dứt lời, tất cả mọi người, kể cả các thành viên đồng thế hệ, lập tức tiến tới đồng thanh
chúc sức khỏe ``cô giáo Y tá'' của họ. Hành động bất ngờ này khiến Choco, vốn thường tự tin với vẻ quyến rũ và thích trêu chọc người khác, đỏ bừng cả mặt vì ngượng ngùng.
Cô lí nhí mấp môi: ``C-Cảm ơn mọi người\dots\ thật là bất ngờ quá.''

Nhìn cô nàng đang đỏ mặt như vậy, Subaru lên giọng châm chọc: ``À-rế? Cô giáo thường
hay trêu học sinh chúng ta thế mà giờ lại ngượng ngùng là sao đây?''

Aqua cũng chen vào: ``Cô không phải xấu hổ đâu ạ! Bọn em chỉ muốn cảm ơn cô thôi
mà, Choco-sensei!''

Nhìn thái độ lúng túng của Sensei, tất cả đều bật cười vì sự đáng yêu hiếm hoi của nữ giáo
viên Y tá này.

Kuon bất chợt vỗ tay, thu hút sự chú ý của cả nhóm. Giọng cậu vang lên dõng dạc: ``Cơ
mà đây không phải là điều anh muốn nói.''

Cả phòng trở nên im lặng hơn, ánh mắt các cô gái dán chặt vào cậu quản lý với vẻ hiếu
kỳ. Cậu con trai tiếp tục: ``Vì hôm nay là `tết Thầy' nên chúng ta sẽ trải nghiệm những gì
mà học sinh, sinh viên hay làm\dots\ hoặc đúng hơn là chúng ta nên làm.''

Cáo trắng Fubuki híp mắt, ngẫm nghĩ trong chốc lát trước khi lên tiếng: ``Đi câu cá sao?''
Rushia chen ngang ngay lập tức sau đó với vẻ tự tin: ``Không, là triệu hồi phép thuật!''

Shion nhướng mày, gõ nhẹ vào tay của nàng Pháp sư: ``Này, cái đó là của chị mà! Đừng ăn
cắp ý tưởng!'' Theo sau đó, những câu trả lời không liên quan cứ liên tiếp được gọi tên,
nhưng rồi Kuon bật cười và xoa tay hết: ``Sai rồi. Mấy cái đó chẳng liên quan gì đến tết
Thầy cả.''

Korone nghiêng đầu, tỏ vẻ ngờ vực, nhưng rồi cũng lên tiếng định tiếp tục đoán: ``Vậy
thì nó là-'' Cậu sau đó ngắt lời, nụ cười toát lên sự ẩn ý: ``\textbf{Khai bút đầu xuân!}''

Cả phòng ``ồ'' lên đầy bất ngờ. Fubuki, vẫn còn thắc mắc, nghiêng đầu hỏi: ``Nó là gì vậy,
Kuon-senpai?''

A-chan nhẹ nhàng đẩy gọng kính lên, và lên tiếng giải thích với sự hiểu biết từ trước:
``Theo tôi hiểu thì nó là tục lệ, hay chính xác hơn là hoạt động học tập để lấy may mắn.''

Kuon gật đầu, tỏ vẻ hài lòng vì có người hiểu ý mình. Tuy nhiên, Mio vẫn tỏ ra bối rối:
``Nhưng mà, cái đó thì liên quan gì đến chúng ta chứ?'' Dẫu rằng cô nàng là một trong
những người được gọi là ``thanh niên nghiêm túc'' của cả công ty, cách nói chuyện của
cậu Tổng khiến cho cô tỏ vẻ dè chừng hơn hẳn.

Kuon nheo mắt, kéo dài câu trả lời để tăng thêm sự hồi hộp: ``Mừng vì em đã hỏi. Để có
thể trải nghiệm đầy đủ nhất, chúng ta sẽ khai bút bằng việc\dots''

Các cô gái đồng thanh hỏi, gần như không thể chờ thêm được nữa: ``Làm gì?''

Cậu chàng thả lỏng vai, thốt lên đến vang cả phòng với nụ cười tinh quái xuất hiện trên
môi: ``\textbf{\dots làm một bài kiểm tra Toán!}''

Bầu không khí đột ngột đóng băng. Cả nhóm sững sờ, những ánh mắt mở to đầy ngạc
nhiên đổ dồn về phía Kuon.

Một vài giây trôi qua trong im lặng, không ai ho he một tiếng trước thông tin đầy bất ngờ
này.

Cốc, cốc, cốc\dots

``Không phải là toán chứ!'' Subaru thốt lên đầy hoảng hốt. Kéo theo đó là đa số mọi
người hoảng loạn với yêu cầu có phần quá quắt với họ.

Miko gần như hét lên, giọng điệu pha chút tuyệt vọng: ``\textbf{Chúng ta\dots\ Chúng ta tàn đời rồi ne!!}'' Mel thì nhẹ nhàng hơn, nhưng cũng chỉ có thể cười trừ: ``C-Có vẻ nước đi này của anh
làm mọi người phát hoảng rồi nhỉ.''

Nhưng không phải ai cũng phản ứng tiêu cực. Pekora khoanh tay, khẽ cười: ``À, cái này
không tệ đâu, peko.''

Okayu gật đầu nhẹ, vẻ mặt trầm tĩnh: ``Xem chừng có vẻ khó đây, nhưng không sao.''
Sora, dù có chút run rẩy, vẫn cố tỏ ra lạc quan: ``C-Chị nghĩ chúng ta sẽ làm được thôi!''

Trong cảm xúc đối lập giữa hoảng loạn và bàng quan về bài kiểm tra môn Toán do Kuon
đề xuất, Ayame lại ghiêng đầu, mắt tròn xoe như thể vừa nghe một khái niệm lạ lẫm:
``\dots Kuon-senpai, toán là gì vậy, yo?''

Aki cũng vậy, vẫn giữ vẻ bình thản, bổ sung thêm với giọng đều đều: ``Lần đầu em được
nghe á.''

Cả nhóm bắt đầu nhốn nháo, mỗi người một câu, khiến không khí trở nên náo nhiệt và
hỗn loạn. Nhìn thấy cảnh tượng ấy, Kuon đành vỗ tay lớn để ổn định trật tự. ``Rồi, rồi
mọi người ơi, bình tĩnh, bình tĩnh lại nào!''

Nhưng quả nhiên, không ai nghe được lời cậu nói. Người tất tả tìm trên mạng để đọc lại
những công thức toán học mà họ đã ``chữ thầy trả thầy'' từ lâu, người thì láo nháo chạy đi
tìm\dots\ sách vở dù họ còn chẳng mang chúng đi.

``Tôi nghĩ mọi thứ diễn ra hoàn hảo đấy, Hikawa-san. Đến bài này tôi cũng không ngờ
tới luôn,'' A-chan nhìn khung cảnh hỗn loạn trước mặt, tỏ vẻ vừa ngạc nhiên nhưng cũng
vừa buồn cười.

Kuon nhếch mép, trả lời với vẻ tinh nghịch: ``Em cũng đã đắn đo suy nghĩ từ trước rồi mà
chị. Chị nên thấy may là em không kiểm tra tiếng Anh đấy.''

Nàng Đeo kính khựng lại một giây, rồi thở dài ngao ngán: ``\dots\ Nhắc mới nhớ, chắc tôi
phải học tiếng Anh lấy bằng TOEIC sớm thôi. Mà cậu nói cũng đúng, thi tiếng Anh thì
có mà tất cả trượt hết.''

\begin{center}
    \uline{Kết thúc hồi tưởng.}
\end{center}

Hiện tại, tất cả các thành viên \hll\ cùng Ryuuhou đang tập trung phía ngoài phòng
314.H1, trường của Kuon. Không khí bên ngoài có phần căng thẳng nhưng lại pha lẫn
tiếng cười nói rì rầm. Trong khi A-chan đang đứng nghiêm túc phổ biến nội quy, bản thân
cậu Tổng đang ngồi bên trong phòng thi, cặm cụi đánh số thứ tự trên bàn với vẻ mặt đầy
kiên định.

Nàng ``Tăng ca'' đẩy gọng kính, cầm tờ danh sách điểm danh và bắt đầu phổ biến quy
định với tông giọng bằng bằng: ``Được rồi, Hikawa-san có viết một vài điều cần làm
trong phòng thi này. Về cơ bản, nó cũng giống quy định ở các trường trung học và đại
học ở Nhật Bản thôi. Nhưng để tránh nhầm lẫn, tôi sẽ đọc cho mọi người nghe.''

Tất cả mọi người im lặng, chờ đợi A-chan bắt đầu phổ biến quy định thi. ``Đầu tiên, không
được mang điện thoại và tài liệu vào phòng thi.''

Chưa để cô nàng tiếp tục, Matsuri chợt ôm đầu, ánh mắt đầy hoang mang: ``Em thấy áp
lực quá\dots\ Y như hồi thi đầu vào trung học vậy\dots''

Subaru cũng hùa theo, giọng điệu nặng nề: ``Tôi cũng\dots\ Chỉ thiếu mỗi giáo viên giám
thị gõ thước kẻ lên bàn nữa thôi.''

Ryuuhou thì bàng quan hơn, vì đây không phải lần đầu cô dự những kỳ thi như vậy. Tuy
thế, cô chỉ thở dài ngao ngán, khi bị anh trai kéo vào tình thế vốn chẳng liên quan gì đến:
``\textit{Tại sao mình lại ở đây chứ\dots}''

Nàng Đeo hắng giọng, tiếp tục đọc quy định trong tờ giấy, không mảy may tới những lời
than thở: ``Thứ hai, chỉ được mang vào phòng giấy thi, bút, máy tính-''

Vừa nghe đến từ ``máy tính,'' Ayame lập tức vác theo\dots\ một chiếc máy tính để bàn (computer)
khổng lồ trước mặt mọi người. Giọng của cô tỏ vẻ ngây thơ: ``Cái này đúng không?''

Mọi ánh mắt lập tức đổ dồn vào nàng Oni. Một vài người cố gắng nhịn cười nhưng không
thành. Subaru thảng thốt: ``G-Gượm đã, Ayame! Bà mang cả máy tính để bàn thật đấy
à!?''

``Thì A-chan bảo được mang máy tính vào mà, yo? Tôi nghe nhầm à?''

Kuon nghe thấy sự náo động từ bên ngoài, ló đầu qua cửa phòng với vẻ mặt khó tin:
``A-Ayame-chan, em định làm bài thi hay là dựng server để chơi game đấy?''

Fubuki cười một cách ranh mãnh: ``Chắc dựng server A*K để chơi với nhau trong giờ thi chứ
gì! Ước gì chị có thể mang bộ máy tính nhà chị sang đây nhỉ\dots''

Thấy mọi người cười để giải tỏa áp lực, A-chan và Kuon chỉ có thể thở dài bất lực: ``Không,
ý tôi là máy tính cầm tay. Calculator ấy, Nakiri-san.''

Ayame đỏ mặt, nhìn chằm chằm vào chiếc máy tính trên tay mình: ``Ồ\dots\ Ra là thế à?
Vậy\dots\ cái này không được à?''

Cả hai giám thị đều cùng đồng thanh: ``Không. Hoàn toàn không được.''

Korone chống cằm, vẻ mặt có chút thất vọng: ``Chán nhỉ. Nếu được mang vào thì chắc thú
vị lắm đây.'' Okayu nhìn sang cô bạn của mình, hùa thêm với vẻ bình thản: ``Nếu được
mang vào chắc ta có thể nói chuyện với nhau nha\~{}''

A-chan thở dài, đáp lại: ``Quy định là vậy rồi, Nekomata-san. Tôi chỉ làm theo yêu cầu
của cậu ta thôi.''

Cả nhóm quay về phía Kuon, người vẫn đang chăm chú đánh số bàn và giả vờ không nghe thấy gì. Cô em gái ở bên ngoài khẽ lẩm bẩm: ``Đúng là không ai chơi khó như vậy bằng
Onii-san cả\dots''

Không khí bên ngoài phòng thi ngày càng trở nên căng thẳng khi A-chan tiếp tục phổ biến
thêm về nội quy. Các thành viên ai nấy đều bồn chồn, vừa lo lắng vừa tò mò trước những
quy định mà Kuon đã đặt ra.

``Tiếp này. Cũng theo bảng nội quy của Hikawa-san, có ba mức cảnh cáo khi bị giám thị
phát hiện vi phạm\dots''

Vừa nghe đến ``cảnh cáo,'' nàng Sói đen liền rụt vai lại, giọng đầy cảnh giác: ``\dots Nghe có
vẻ nghiêm trọng đây.''

``Đúng đấy. Lần đầu bị nhắc nhở, bài thi sẽ bị trừ 25\% số điểm.''

Mel, người vốn luôn điềm tĩnh, lần này cũng không giấu được sự ngạc nhiên: ``Sao nghe
có vẻ nặng nề thế\dots?''

Nàng Đeo kính vẫn tiếp tục phổ biến với giọng bàng quan: ``Chưa nặng bằng hai cái sau
đâu. Lần thứ hai sẽ bị trừ nửa số điểm bài thi. Và nếu vi phạm lần ba\dots''

Roboco, đoán trước được câu trả lời, nghiêng đầu hỏi trong sự lo lắng: ``Em nghĩ là sẽ bị
đuổi khỏi phòng ạ?''

A-chan gật đầu: ``Đúng rồi. Đồng nghĩa là em sẽ bị đình chỉ thi ngay lập tức và chịu phạt
thêm nữa.''

Mio trợn tròn mắt, giọng hoang mang: ``Có cả phạt nữa à!? Chỉ là một hoạt động ngày
Tết thôi mà anh ấy làm gắt thế\dots''

A-chan khẽ nhún vai, chỉ vào tờ giấy trên tay mình: ``Trong này có quy định thưởng phạt
rõ ràng mà. Có ý kiến gì tí nữa hỏi cậu ta ấy.''

Ryuuhou như được dịp, lên tiếng an ủi mọi người: ``Lúc bọn em thi cũng nghiêm túc như
vậy đấy, các chị à. Không đùa được đâu.''

Noel, một tay chống cằm suy nghĩ, một tay khẽ vỗ vai đồng đội bên cạnh: ``Không biết
liệu anh ấy có thể nới lỏng cho ta không nhỉ\dots''

Pekora cảm thấy lạnh gáy, lẩm bẩm với sự lo ngại: ``Thật. Tôi nghe đến đây đã thấy rén
rồi đó, peko.''

Thấy mọi người tỏ vẻ băn khoăn, nàng Đeo kính, người hiểu khá rõ tính cách của Kuon,
nhấn mạnh: ``Tôi e là không. Hikawa-san được biết đến là một trong những nhân viên
nghiêm túc nhất công ty. Không đời nào cậu ấy sẽ cho qua đâu.''

Tất cả đồng loạt hít một hơi dài, mắt nhìn nhau đầy sợ sệt. ``T-Thật à\dots'' Ryuuhou chỉ có
thể mỉm cười, tỏ vẻ đông cảm: ``Onii-san thực sự làm rất gắt gao với những vấn đề này
đấy ạ\dots\ Em nghĩ chúng ta sẽ làm được thôi!''

A-chan cố trấn an, vỗ tay để thu hút sự chú ý: ``Cứ nghiêm túc thi là được. Bây giờ tôi sẽ
nói nốt về phần thưởng phạt và cơ cấu đề thi đây.''

Không khí ngoài hành lang vừa hồi hộp, vừa có chút nặng nề. Ai cũng cảm thấy như mình
đang đứng trước một kỳ thi thật sự, chứ không phải chỉ là một hoạt động trải nghiệm Tết.

``Khoan đã, A-chan!'' Ayame dừng lại, rồi tỏ vẻ băn khoăn: ``Chị có nói rõ hơn được vẻ
quy chế thưởng phạt được không?''

``Tôi đang định nói đây. Cứ bình tĩnh để tôi nói.'' Chưa để A-chan tiếp tục, Fubuki đã leo
tới, đọc những gì được viết trên tờ nội quy đó: ``Về cơ bản, có bốn mức thưởng phạt trong
bài thi này. Ai từ 8.25 đến 10 điểm sẽ được lì xì mười nghìn yên!''

Nghe đến giải thưởng rất lớn kia, tất cả mọi người đều không giấu khỏi sự mừng rỡ. Tuy
vậy, Sora tỏ vẻ ngạc nhiên: ``8,25 đến 10? Tại sao điểm thấp vậy?''

A-chan giải thích: ``Tôi nghĩ là tính theo thang 10, không phải thang 100. Etsunan tính
điểm theo thang 10 mà.'' Cô liếc sang tờ giấy đó, và nhấn mạnh: ``Đặc biệt là, Hikawa-san
không chấp nhận làm tròn điểm. Tức là nếu ai đạt 8.25 là được 8.25 luôn, không làm tròn lên, cũng không bị làm tròn xuống.''

Tất cả mọi người đều gật gù hiểu. Matsuri lên tiếng: ``Thật á!? Vậy từ 5 đến 8 điểm thì
sao?''

Cáo trắng tiếp tục: ``Nghe nói là anh ấy sẽ thưởng năm nghìn yên.'' Bất chợt, mặt cô nàng
tỏ vẻ nghiêm trọng hẳn: ``Nhưng nếu ai đó tron các cô chỉ được từ 3 đến 4.75 thì\dots''

``Thì sao?'' Shion tỏ vẻ ngạc nhiên.

``Thì phải chịu phạt làm theo yêu cầu của những người điểm cao nhất!'' Fubuki thốt lên,
tỏ vẻ hào hứng. Đến đây, tất cả mọi người đều tỏ vẻ hoang mang.

``Làm theo yêu cầu? Kiểu như làm gì cơ\dots?'' Mel lên tiếng, bắt đầu đổ mồ hôi hột.

Nàng Cáo trắng cười gian: ``Anh ấy không ghi rõ ở đây, nhưng nghe nói là họ sẽ tự quyết
định. Chắc chắn là sẽ thú vị lắm\~{}''

Subaru lên tiếng, hỏi với giọng sợ sệt: ``Vậy\dots\ từ ba điểm trở xuống thì sao?'' Mio nghe
vậy, thở dài: ``Chắc chắn là sẽ có hình phạt gì đó không dễ chịu gì cho cam đây.''

A-chan gật gù, dù cô có hơi khẽ rùng mình vì hình phạt có phần\dots\ ác liệt: ``Ookami-san
nói đúng đấy. Những người điểm dưới 3 sẽ được `vinh dự' đeo bảng `\textbf{Tôi dốt toán}' trước ngực và đi diễu hành vòng quanh trường.''

Đa số mọi người trố mắt ngạc nhiên với hình phạt có phần khắc nghiệt đó. Subaru trố
mắt: ``C-Cái gì!? Đùa chứ, sao nghe giống trừng phạt công khai thế!''

Đến cả cô em gái của Kuon cũng không tin vào những gì mà cô vừa nghe: ``E-Em cũng
không nghĩ Onii-san sẽ nghĩ ra hình phạt như vậy đấy\dots''

Fubuki bật cười, nói với giọng ma mãnh: ``Không đâu, hình như cái này là ý tưởng của
Shion-chan và Noel-chan đấy!''

Tất cả mọi người quay về phía nàng Phù thủy và nàng Hiệp sĩ với những ánh mắt hình
viên đạn.

\begin{dialogue}
    \speak{Subaru} Thế này là thế nào hả?
    \speak{Mio} Làm sao mà hai người có thể nghĩ ra được hình phạt oái oăm như vậy chứ?
    \speak{Shion} Có sao đâu? Nếu mọi người không muốn bị phạt thì thi nghiêm túc vào đi chứ.
    \speak{Noel} Ờ thì\dots\ chỉ là bọn tớ muốn tạo chút động lực cho mọi người thôi mà.
\end{dialogue}

Korone thốt lên, giả vờ tỏ vẻ nghiêm túc: ``Cơ mà hai người nói thế lại càn khiến người
ta trốn thi luôn đấy chứ.''

Dẫu cho hai người họ giải thích thế nào, họ vẫn phải chấp nhận với hình phạt này. Nàng
Sói đen tiếp tục hỏi: ``Thế còn cơ cấu đề thi thì sao? Có phải là mấy bài viết luận dài
ngoằng không?''

A-chan bình tĩnh trả lời, tay vẫn đang cầm tờ giấy nội quy: ``Không phải đâu. Đây là một
bài trắc nghiệm năm mươi câu, mỗi câu có bốn đáp án.''

Câu trả lời này khiến tất cả thở phào nhẹ nhõm, vì ít nhất họ vẫn còn có thể làm được bài.

``Ít nhất thì chúng ta không phải nộp giấy trắng\dots'' Matsuri thốt lên, lau mồ hôi trên trán.

``Cơ mà tỷ lệ đoán trúng cũng không cao lắm đâu,'' Mio nói, tỏ vẻ lo ngại hơn hẳn, ``một
câu trả lời có 4 đáp án, chỉ có 25\% là đúng thôi.''

``Tôi cảm giác là đoán đúng cả 50 câu có khi còn khó hơn cả trúng xổ số đấy chứ, yo,''
nàng Oni khoanh tay lại, miệng vẫn nở nụ cười gian xảo. ``Tỷ lệ là \(\frac{1}{4^{50}}\) á.''

Câu trả lời của cô càng khiến mọi người lo lắng hơn, nhất là trong bối cảnh mà gần như
mọi người không còn nhớ được khái niệm gì về môn toán cả.

Miko lên tiếng: ``Vậy chúng ta sẽ trả lời như thế nào, ne?''

A-chan giải thích: ``Mỗi người sẽ phát đề thi và một phiếu trả lời trắc nghiệm. Trên phiếu
trả lời đó, mọi người sẽ phải ghi mã đề và số báo danh lên phiếu trả lời. Tô cẩn thận là
được.''

``Vậy thì thời gian làm bài thì sao, ne?''

A-chan lật mặt sau của tờ nội quy, rồi khẽ lắc đầu: ``Trong này không có ghi. Chắc
Hikawa-san sẽ giải thích rõ hơn thôi.''

Ngay sau khi cô dứt lời, Kuon ra hiệu lệnh cho A-chan bắt đầu đọc danh sách thí sinh
phòng thi, được đánh số từ 01 đến 19. Các chỗ ngồi được sắp xếp theo bảng chữ cái thay
vì theo thế hệ.

A-chan bắt đầu giở tờ thứ hai chứa danh sách những người đến thi một cách to và rõ ràng.
``Aki Rosenthal-san?''

``Có!''

``Minato Aqua-san?''

Nghe đến tên, nàng Hầu gái lúng túng giơ tay: ``C-Có ạ!'' rồi vội vã vào trong phòng.

Cứ như vậy, nàng Đeo kính tiếp tục đọc tên lần lượt từng người, cho đến người cuối cùng.
``Murasaki Shion-san?''

Nàng Phù thủy giơ tay, tỏ vẻ tự tin: ``Có!''

``Hikawa Ryuuhou-san?''

``Có ạ, thưa chị.''

Từng người lần lượt bước vào phòng và ngồi theo số thứ tự được ghi trên bàn.

Sau khi mọi người ổn định chỗ ngồi, Kuon bước lên phía bảng với xấp đề thi trên tay.
Chưa để ai nói một lời gì, cậu nói: ``Bài thi Toán của chúng ta sẽ kéo dài 120 phút. Được
phép sử dụng máy tính bỏ túi. Sẽ có ba cán bộ coi thi, ai vi phạm thì theo quy định mà xử
lý.''

``Ba người coi thi ạ?'' nàng Sói đen tỏ vẻ thắc mắc, ``Nhưng mà bọn em chỉ thấy anh với
A-chan đang đứng đây thôi mà?''

Sora nhìn các bàn xung quanh, rồi chợt nhớ ra: ``Mà, Choco-sensei đâu rồi nhỉ? Chị nhớ
là cô cũng đi theo chúng ta vào đây mà\dots''

Rushia cũng cảm thấy như vậy: ``Đúng á-nanodesu! Cô ấy trốn đi đâu rồi không biết\dots''

Vừa lúc đó, cánh cửa phòng bật mở. Một dáng người quen thuộc xuất hiện với chiếc bảng
tên ``Giám thị Choco'' đeo trên cổ.

``Xin lỗi mọi người nha, cô đến muộn\~{} Nhưng như Kuon-sama nói đấy, cô sẽ là giám thị
thứ ba của buổi thi này.'' nàng Y tá ung dung bước vào, nói với giọng trầm nhưng quyến
rũ.

Tất cả theo dõi cô nàng ấy với vẻ sửng sốt. Subaru há hốc miệng, ngạc nhiên: ``\textbf{Ể-Ể!? Choco-sensei là giám thị ạ!?}''

Mio thì kinh ngạc, nhưng rồi cũng hiểu rõ được vấn đề: ``Chẳng trách từ nãy giờ chẳng thấy cô ấy đâu\dots''

Thấy mọi người lại nhốn nháo do bất ngờ, A-chan - giám thị thứ hai của buổi thi - bình
tĩnh giải thích: ``À, cô ấy được giao vai trò này vào phút cuối. Hikawa-san là người quyết
định mà.''

Lời giải thích của cô không khiến mọi người không khỏi ngạc nhiên, thậm chí còn làm
cho cả phòng nhốn nháo hơn. Kuon thở dài, đập mạnh vào bàn và nói: ``Rồi, lớp ơi, lớp
ơi, trật tự nào. Còn ồn ào nữa là anh cho mọi người không điểm hết đấy.''

Cả phòng lập tức yên lặng, không khí nghiêm túc bao trùm. ``Được rồi, anh đến đâu rồi
nhỉ\dots''

Cậu Tổng cầm xấp đề thi lên, phát đề cho mọi người. ``Đề anh sẽ để úp, khi nào có tín
hiệu bắt đầu làm bài thì mới lật lên trên.''

``Để làm gì vậy, ne?'' Miko nhíu mày. Kuon đáp: ``Để tránh gian lận đấy mà.'' Cậu đi
xuống, phát mỗi người một đề thi, cùng lúc với A-chan phát phiếu trả lời trắc nghiệm.

Trong lúc phát đề, cậu tiếp tục dặn dò: ``Tuyệt đối không sử dụng tài liệu, không dùng
điện thoại, không trao đổi bài. Ai nhìn thấy có điện thoại trong người, dù sử dụng hay
không, lập tức 0 điểm tại chỗ và chịu hình phạt.''

Câu nói khiến mọi người rén hơn, không ai dám ho he một tiếng, nhất là khi mỗi người
ngồi một bàn riêng biệt và cách nhau rất xa. ``Còn điện thoại hay là gì thì mang hết lên
bàn giáo viên nha.''

Lập tức, tất cả mọi người hối hả đi đến bàn giáo viên, nơi mà Kuon sẽ ngồi trông coi, và
đặt điện thoại của họ trên đó, trước khi hối hả quay về bàn.

Cậu quay trở về vị trí trông thi của mình, hỏi: ``Có ai còn ý kiến gì không?''

Gần như tất cả mọi người trong phòng im re, không có lấy một tiếng nào. Khi cậu Tổng
trước khi nói thêm một câu nào, Mel giơ tay yếu ớt - dường như vì sức nặng tại căn phòng
này: ``Bọn em có được phép đi vệ sinh không?''

Kuon gật đầu, nói với sự nghiêm túc, dù cố gắng làm bầu không khí dịu hơn đôi chút:
``Đương nhiên rồi.'' Tất cả mọi người đều thở phào nhẹ nhõm.

``Nhưng chỉ giới hạn một người đi và không được đi quá ba phút,'' cậu nói thêm.

Subaru nghe vậy, tỏ vẻ tò mò: ``Nghĩa là bọn em không phải nộp bài để đi vệ sinh ạ?''

Cậu đáp: ``Không. Ở đây khác so với ở Nhật, các em có thể đi giải quyết nỗi buồn thoải
mái, nhưng không được để người khác nhìn bài.''

Cô nàng rón rén hỏi tiếp: ``Vậy đề toán này của anh có khó không?'' Cậu mỉm cười đầy
ẩn ý, đáp lại: ``Dễ í mà. Mấy đứa còn có máy tính nữa cơ, sợ gì.''

Đến đây, không còn ai có thắc mắc gì với cậu nữa. Tất cả mọi người đã an tâm phần nào
khi họ có thể làm bài thi thoải mái hơn so với hồi ở quê nhà.

``Bây giờ là tám giờ mười phút. Ta sẽ làm đến mười giờ mười thì nộp bài. Ai hết giờ
không nộp, coi như không điểm.''

Kuon liếc nhìn chiếc đồng hồ, rồi dõng dạc nói trước mười chín người:

\begin{center}
    ``Thời gian làm bài của các em bắt đầu.''
\end{center}

Ngay khi có hiệu lệnh, tiếng soàn soạt của tờ giấy vang lên, cùng lúc với những tiếng bút
chì xoèn xoẹt đầu tiên hòa vào nó. Cả phòng bây giờ chỉ có những tiếng của giấy và bút,
hoàn toàn không có một tiếng động nào khác.

\ct

Đúng như cậu Tổng nói, đề toán mà cậu chuẩn bị cho mọi người đều nằm ở những kiến
thức sơ đẳng, nhưng\dots\ có lẽ chỉ ``sơ đẳng'' với những người rất thích học toán như cậu.

``\textit{Trong này còn có cả phân số này\dots}'' Mio nghiêng đầu nhìn đề thi, cẩn thận nhẩm từng
dòng. Cô đặt bút xuống, bắt đầu viết vài phép tính thử nghiệm, nhưng ánh mắt không
giấu được chút căng thẳng.

Ở một bàn khác, Korone gãi đầu, nhíu mày với tờ đề thi của mình: ``\textit{Còn có hằng đẳng
    thức nữa\dots\ Mấy cái bình phương với căn thức trông không quá khó, nhưng mà\dots}'' Giọng
cô nhỏ dần, ánh mắt đăm chiêu khi nhìn thấy một bài toán về phương trình bậc hai.

Nhưng với Pekora, cô nàng khẽ thốt lên khi nhìn tờ đề của mình: ``Sao còn có cả chia đa
thức thế này, peko!?'' Cô vừa nói vừa chỉ vào một bài toán dài ngoằng, trông như một cơn
ác mộng.

DĨ nhiên, đề thi của ba người họ chưa là gì so với một nàng Vu nữ, khi cô vừa nhìn qua
các câu phía sau của đề\dots

``Đề thi so khó quá vậy, ne\dots'' mặt cô tái mét hẳn, khẽ rên rỉ. Ngay sau đó, cô nàng gục
xuống bàn, ngất lịm. Sora ngồi trước cô thấy vậy hoảng hốt, lay nhẹ vai của Miko, cố
gắng gọi: ``Miko-chan, dậy đi, Miko-chan!''

Mel ngồi gần đó, nhìn nàng Vu nữ đang xỉu đó, khẽ nghiêng người nhìn vào đề thi của cô
ấy: ``Không biết đề của chị ấy trông thế nào nhỉ\dots''

Nàng Dơi cầm lên tờ đề, và ngay lập tức tá hỏa khi thấy các câu về tổ hợp, chỉnh hợp,
và cả\dots\ định thức ma trận. Đề thi của chính cô nàng thậm chí còn có bài toán về xâu và
automata, thứ khiến cô cảm giác như đang đọc một ngôn ngữ ngoài hành tinh.

Ở bên góc phòng, Fubuki cũng không khá khẩm hơn, thậm chí có thể nói là ``khoai khoẳm''
hơn, khi cô run rẩy nhìn những bài toán hóc búa: ``Đề của mình còn có tích phân, đạo
hàm, đại số tuyến tính và mã hóa\dots''

Vừa dứt lời, cô nàng cũng ngất xỉu ngay trên bàn, khiến Mio và Ayame phải luống cuống
lay gọi mãi mà không thấy động tĩnh.

Trong lúc mọi người còn đang bấn loạn vì đề thi oái oăm đó, chỉ có duy nhất là Ryuuhou
là sát với những gì cô bé đã được học ở trường - với một học sinh năm ba trung học, các
câu hỏi đều vừa sức và dễ dàng. Cô gái trẻ mỉm cười, bình tĩnh đặt bút làm bài, khiến các
thành viên khác càng thêm ghen tị, nhưng không dám lên tiếng thắc mắc.

Từ phía trên cao, Kuon cười nhạt khẽ quan sát phản ứng hoảng loạn của mọi người: ``Đấy,
các em cứ làm từ từ đi. Nếu khó quá thì khoanh bừa cũng được. Nhớ đừng quên luật chịu
phạt nếu dưới 3 điểm nhé.''

Câu nói của cậu càng khiến không khí phòng thi thêm phần căng thẳng.

A-chan và Choco sau khi đọc xong một trong những đề của cậu con trai cũng không khỏi
ngỡ ngàng:

\begin{dialogue}
    \speak{A-chan} Sao đề của cậu khó thế này, Hikawa-san!?
    \speak{Choco} Đến em còn thấy choáng nữa đây\dots\ Thật đấy!
\end{dialogue}

Kuon chỉ nhún vai, khẽ lắc đầu: ``Toàn kiến thức cơ bản đấy hai người ạ. Em nghĩ mọi
người cũng từng học qua hoặc ít nhất là thấy sơ qua rồi mà.''

Dứt lời, một số thành viên ở dưới đó liền lên tiếng bất bình. Subaru là người phản đối
mạnh mẽ nhất: ``Cơ bản cái con khỉ! Ai mà làm được mấy cái đó được chứ!''

Rushia cũng tiếp lời: ``Đúng đấy! Bọn em chỉ làm nổi phép toán kiểu cộng, trừ, nhân,
chia của học sinh tiểu học thôi!''

Ayame cầm tờ đề lật qua lật lại, nét mặt đầy bối rối: ``Em thấy đề này\dots\ nhìn hơi bất khả
thi thật.''

Kuon bật cười, một tay chỉ vào đề, một tay khoanh trước ngực: ``Thì anh bảo đây là `\uline{kiến thức sơ đẳng}', chứ có phải `\uline{kiến thức tiểu học}' đâu.''

Nàng Pháp sư nghe vậy, liền tỏ vẻ kịch liệt lập luận: ``Nhưng nếu như thế thì chẳng phải
bọn em trượt sạch hay sao!?''

Trái ngược với những lời ca thán, Roboco, với vẻ tự tin thường thấy, khẽ nhếch môi sau
khi đọc qua đề: ``Ừm, em thì có thể tính toán được nhanh lắm đó.''

Lập tức, nàng Vịt phản ứng mạnh, vung tay lên trời: ``Ê, nhưng chị là người máy! So
sánh làm gì! Chơi vậy ai chơi lại!?''

Nghe những lời than thở từ họ, giám thị chính chỉ thở dài, gõ nhẹ lên mặt bàn vài lần để
thu hút sự chú ý: ``Đấy là việc của mấy đứa thôi. Câu nào không làm được thì tạm bỏ qua.
Câu nào dễ thì làm trước. Anh tưởng mọi người phải biết quy luật này rồi chứ?''

Thấy cả lớp vẫn đang ồn ã, người tiếp tục phản đối, người cố gắng gọi đồng đội tỉnh dậy,
cậu tiếp lời: ``Với lại, đây là bài trắc nghiệm, mỗi câu có bốn đáp án, 25\% cơ hội đúng.
Không biết thì cứ khoanh đại. Coi như chơi xổ số ấy.''

Noel gãi gãi cằm, có vẻ không tin tưởng lắm: ``Khoanh đại\dots\ Kiểu đó cũng được à?'' Cậu
gật đầu, đưa ra một ví dụ mà nghe qua đã thấy đầy hên xui may rủi: ``Ừ. Ví dụ, em cứ \href{https://www.facebook.com/watch/?v=445235674120478}{chọn toàn C} chẳng hạn. Đáp án C thường là câu bẫy, nhưng ít nhất thì\dots\ nó có cơ hội.''

Ngay khi cậu đưa ra ví dụ ``có cũng như không,'' Sora liền cau mày: ``Nhưng nếu làm thế
thì hỏng hết mục đích bài kiểm tra còn gì, Kuon-senpai.''

Kuon nở một nụ cười đầy bí hiểm, nhún vai đáp: ``Nói thế này thì có lẽ mấy đứa không
tin, nhưng bọn anh toàn làm thế suốt hồi học đại học. Bản chất bài kiểm tra là để kiểm
tra\dots\ những gì mà mình không biết thôi.''

Câu nói này khiến tất cả mọi người trong phòng chợt cứng họng. Không ai biết nên phản
bác hay đồng tình, chỉ còn lại những giọt mồ hôi lạnh lăn trên trán. Trong đầu của mọi
người giờ đây cũng không biết phải thắc mắc gì hơn, ngoài đúng một ý nghĩ: ``\textit{\dots Không biết anh ấy đã học kiểu gì nữa\dots}''

``Thôi, mấy đứa cứ tập trung làm bài đi. Học tài thi phận thôi.''

Sau khi Kuon tuyên bố bắt đầu làm bài, không khí trong phòng thi trở nên yên ắng đến lạ
thường. Tiếng lật giấy loạt xoạt và tiếng bút chì kẽ trên phiếu trả lời trắc nghiệm là những âm thanh duy nhất vang lên. Dù vậy, sự căng thẳng hiện rõ trên từng khuôn mặt.

Fubuki và Miko, sau màn ngất xỉu bất đắc dĩ, được mọi người cố gắng lay gọi. Đến khi
tỉnh lại, cả hai vẫn trong trạng thái hoảng loạn:

\begin{dialogue}
    \speak{Miko} Tại sao đề toán này khó thế, ne!?
    \speak{Fubuki} Em thậm chí còn không nhớ lần cuối mình làm bài toán là khi nào nữa\dots
\end{dialogue}

Sora nắm lấy tay Miko, nói nhỏ, giọng an ủi: ``\hl{Không sao đâu, làm hết sức mình là đượ.
    Mình tin là cậu có thể làm được, Miko-chan!}''

Mel đứng bên cạnh cũng gật đầu, nhẹ nhàng động viên: ``Đúng đấy! Với lại, đây là bài
trắc nghiệm mà, cứ chọn bừa cũng được!''

Những lời an ủi của họ cuối cùng cũng giúp Miko và Fubuki trấn tĩnh hơn một chút, đủ
để bắt đầu cầm bút và thử làm bài. Dù vậy, cứ vài phút, tiếng than thở của nàng Vu nữ
lại vang lên: ``\textit{Không hiểu sao bài này lại phải tính cả cái tích phân lằng nhằng này chứ
    ne\dots\ Mấy cái này thì giúp gì được trong đời thực cơ chứ, ne!}''

Ở phía đầu lớp, Kuon, Choco và A-chan đã chuẩn bị sẵn máy quay để ghi lại khoảnh khắc
đặc biệt này. Nàng Đeo kính cầm máy quay, còn cậu Tổng đứng bên cạnh, vừa nhìn vừa
chỉ tay, không quên thêm lời nhận xét: ``Chị thấy nét mặt của Subaru-chan không? Nhìn
có khác gì đang gặp ác mộng đâu.''

A-chan khẽ bật cười, hướng máy quay về phía nàng Vịt, người đang vò đầu bứt tai khi
cố gắng giải một bài toán tổ hợp: ``Chị nghĩ phải quay lại cảnh này. Mấy đứa về xem lại
chắc sẽ tự cười bản thân đấy.''

Choco cũng bật cười khi nhìn vẻ mặt lúng túng của tất cả mọi người: ``Đề thi mà anh cho
đúng là dị thật đấy\dots\ Không biết họ sẽ xử trí thế nào nhỉ.''

Korone ngồi cách đó không xa, đôi tai chó của cô khẽ cụp xuống, vừa nhai bút vừa nhìn
chằm chằm vào câu hỏi: ``\textit{Hmm\dots\ Nếu chọn sai thì chọn đáp án nào nhỉ? A hay B nhỉ\dots?
    Hay cứ C luôn cho chắc?}''

Noel ở phía bên kia phòng thì đập tay lên bàn khe khẽ, miệng lẩm bẩm: ``\textit{Hình như mình
    nhớ công thức này\dots\ Không đúng, sai rồi! Công thức nào mới đúng nhỉ!?}''

Cả phòng thi lúc này, dù im lặng, nhưng ai cũng toát lên vẻ căng thẳng. Chỉ có Kuon,
đứng một góc, vừa liếc nhìn đồng hồ đeo tay của mình, vừa mỉm cười như thể rất hài lòng
với bài kiểm tra đầy thách thức của mình.

\ct

Cũng có những tình huống buồn cười khi cả ba người tuần tra trong phòng thi. Mỗi lần
đi qua một bàn, họ lại nhận được những ánh mắt cầu cứu hoặc sự hoảng loạn từ các thí
sinh.

Đến bàn của Subaru, cậu con trai thấy nàng Vịt đang cắn bút, nhìn chằm chằm vào hai
đáp án B vừa tô. Khi thấy bóng cậu, cô lập tức giơ tay lên hỏi: ``Kuon-senpai, hai câu liên
tiếp em chọn B rồi\dots\ không biết em sai ở đâu nhỉ?''

Cậu con trai nhướng mày, cười khì: ``Sao lại hoảng hốt thế? Cứ tin vào bản thân đi chứ.''

``Nhưng mà hai câu giống nhau thì em thấy cứ sai sai! Em cũng đã từng bị một lần như
thế từ trước rồi\footnote{Cụ thể hơn, trong một buổi stream vào năm 2021, \href{https://www.youtube.com/watch?v=4gsPSOFntgk}{Subaru đã từng thực hiện phép đặt tính rồi tính, với hai câu đầu cho kết quả giống nhau.}}!''

Chưa kịp trả lời nàng Vịt, từ bàn bên cạnh, Noel bỗng bật ra một tiếng kêu khẽ: ``Chết
thật! Sao bốn câu liền mình chọn A thế này!?''

Kuon quay lại nhìn cô nàng đang hoảng loạn, khóe miệng nhếch lên: ``Thế còn nhẹ chán
ấy. Bọn anh còn chọn liên tiếp một đáp án nguyên cả đề ấy chứ.''

Lời nói của cậu càng khiến Subaru xanh mặt hơn, rồi cô cúi gằm xuống, tập trung vào bài
làm của mình. A-chan đứng sau lưng, khoanh tay hỏi: ``Cậu đang nói ra bí mật đấy à?''

Cậu cười trừ, đáp: ``Cái này em nói thật. Thật sự em đã từng làm bài thi còn mất não hơn
như thế.''

Mio nghe cuộc hội thoại từ ba người, nhìn lên với ánh mắt hoài nghi: ``Này, Kuon-senpai,
sao em thấy giống như anh đang cố thao túng tâm lý bọn em nhỉ?''

Cậu Tổng quay sang nàng Sói đen, khẽ nhún vai: ``Anh nghĩ là mấy đứa đừng nên nghĩ
quá nhiều. Với cả tập trung làm bài đi, còn tận chín mươi phút nữa cơ mà.''

Choco, đứng ở góc phòng, chỉ im lặng quan sát. Thấy ánh mắt của nàng Y tá, Kuon ghé
lại gần, hỏi nhỏ: ``Em có vẻ trầm tư nhỉ, Sensei.''

Choco nhìn vẻ mặt căng thẳng của họ, khẽ cười mỉm: ``Mọi người đúng là dễ thương
thật.''

``Ừ, thì dễ thương đấy,'' cậu đảo mắt, liếc nhìn về chiếc đồng hồ đeo tay, ``nhưng bài thì
vẫn đang làm dở. Đi thôi, còn nhiều đứa nữa cần anh kiểm tra.''

Cả ba tiếp tục đi tuần, để lại Subaru và Noel trong nỗi lo lắng về chuỗi đáp án của mình.

\ct

Trong một góc phòng, Fubuki đang xoay xở với bài thi của mình. Máy tính của cô bỗng
nhiên dở chứng, không lên nguồn dù cô đã thử nhấn mọi nút. Buộc phải làm toán bằng tay,
nàng Cáo trắng cau mày nhìn câu hỏi có vẻ ``quá sức'' của học sinh cấp hai: \(3 \times 6 + 4 \times 2\).

``Ê-tô\dots\ 3 nhân 6 là\dots\ 18\dots''

Kuon vừa đi tuần đến gần, nhìn vào bài làm của cô, khẽ gật đầu: ``Ừm-hứm. Tiếp đi.''

Fubuki nhíu mày, căng thẳng hơn. ``Rồi cộng với 4\dots\ ra\dots\ 22!''

Cậu con trai che miệng, cố nhịn cười trước sự tính toán ``tài tình'' của cô. Nhưng trước
khi kịp nhắc, cô nàng đã tiếp tục: ``Rồi nhân với 2 nữa là\dots\ 44! Chắc chắn là nó rồi!''

Cậu nhìn nàng Cáo trắng với ánh mắt nửa buồn cười, nửa bất lực: ``Wow. Bà hoàng tính
toán có khác.''

Fubuki đẩy ngực tự hào, ánh mắt sáng lên: ``Chuyện! Phép toán này là chuyện nhỏ! Mấy
cái này mà khó với Fubuki-sama à?''

Ngay từ phía cậu, A-chan cũng không nhịn được cười: ``Hikawa-san, cậu không định sửa
cho cô ấy à?''

Cậu nhún vai: ``Giám thị chúng ta có được phép nhắc đâu. Với lại, để em ấy tự phát hiện
ra cũng hay mà.''

Thấy rằng hai người họ đang nhắc tới cô, Cáo trắng nghiêng đầu nhìn hai người: ``Sao
thế? Có gì không đúng à?''

Kuon nghe thế, vội xua tay: ``Không, không, cứ tiếp tục đi. Anh xem thử người khác làm
bài đây.''

Cậu con trai và A-chan tiếp tục di chuyển, để lại Fubuki loay hoay với bài thi. ``Tính lại
cho chắc vậy\dots\ 18 cộng với 4 rồi nhân với 2\dots\ đúng rồi mà!''

Cô nàng gật đầu tự mãn, lại tiếp tục với câu tiếp theo, mặc kệ ánh mắt thương cảm từ
Choco đang đứng gần đó.

\ct

Ở một góc phòng khác, Rushia đang chăm chú làm bài, ánh mắt dán chặt vào một câu
phương trình tưởng chừng đơn giản: \(\log_2\left(\frac{7}{4} + 2\right) - \log_2\left(\frac{5}{7} - 8\right)\). Sau một hồi tính toán, cô đã tự tin chọn đáp án và chuẩn bị chuyển sang câu tiếp theo.

Vừa lúc đó, Kuon đi ngang qua. Nhìn thấy đáp án mà cô vừa chọn, cậu không nói gì mà
chỉ khẽ lắc đầu, rồi tiếp tục bước đi.

Nhìn thấy cái lắc đầu của cậu, nàng Pháp sư giật mình, ánh mắt lo lắng nhìn theo bóng
lưng của cậu: ``\textit{C-Chẳng nhẽ mình làm sai ở đâu sao!?}''

Cô nàng bắt đầu nghi ngờ bản thân, nhìn lại từng bước tính toán: ``\textit{Có thể là ở phép trừ
    logarit này\dots? Hay là phép cộng\dots?}''

Cô tái mặt, cẩn thận xóa đáp án vừa chọn, bắt đầu tính lại từ đầu. Nhưng càng tính, cô
càng cảm thấy rối trí, đầu óc như quay cuồng giữa các con số. ``\textit{Nhỡ đâu mình làm sai
    thật thì sao\dots\ Nhưng nếu mình sửa lại rồi vẫn sai thì sao!? Aaaaa!}''

A-chan đứng phía xa quan sát, thấy vẻ mặt hoảng loạn của nàng Pháp sư, liền khẽ huých
tay Kuon: ``Này, cậu làm gì mà khiến cô ấy hoảng loạn thế?''

Kuon nhún vai, tỏ vẻ vô tội: ``Có làm gì đâu. Em chỉ lắc đầu thôi mà.'' Nàng Đeo kính
nhìn cậu chằm chằm, thở dài bất lực: ``Đúng là kiểu đánh tâm lý của cậu\dots''

Trong khi đó, Rushia vẫn đang loay hoay với đáp án của mình, cuối cùng lại quay lại đáp
án ban đầu mà nàng đã chọn: ``Chắc là cái này đúng\dots\ Hy vọng là đúng\dots''

Cậu con trai nhìn cô nàng thở phào nhẹ nhõm, khẽ mỉm cười: ``Xem kìa, mèo vẫn hoàn
mèo. \hl{Đáp án em ấy chọn đúng rồi đấy.}''

Choco bật cười khẽ, thì thầm với A-chan: ``Anh ấy đáng sợ thật đấy. Không cần làm gì
cũng khiến người ta tự nghi ngờ bản thân.'' Nàng Đeo kính cũng chỉ thở dài, ngầm công
nhận với ``tài'' làm giám thị của cậu.

\ct

Aqua với vẻ mặt căng thẳng, dán mắt vào một câu hỏi xác suất trên tờ đề, tay cô liên tục
viết rồi xóa trên giấy nháp. Đề bài trông có vẻ đơn giản: ``Một túi có 5 viên bi, gồm 3
viên đỏ và 2 viên xanh. Lấy ngẫu nhiên 2 viên bi, xác suất để cả hai viên bi đều là đỏ là
bao nhiêu? Lấy giá trị gần đúng.''

Nàng Hầu gái hoảng loạn, đôi mắt cô xoay như chong chóng: ``\textit{Xác suất\dots\ là gì nhỉ? Trời ơi, mình đâu có nhớ cái này đâu\dots}''

Cô hít một hơi thật sâu, bắt đầu lấy máy tính toán ra để tính: ``\textit{Đầu tiên, số cách chọn 2 viên bi từ 5 viên là\dots\ \(C^2_5\) hay \(\binom{2}{5}\)\dots}'' và tiếp tục hí hoáy: ``\textit{nhân 4 rồi chia 2\dots\ là 10!
    Đúng rồi!}''

Sau đó, cô tiếp tục: ``\textit{Số cách chọn 2 viên đỏ từ 3 viên là\dots\ \(C^2_3\) hay \(\binom{2}{3}\)\dots\ 3 nhân 2 chia 2\dots\ ra 3!}''

Ghi chép đầy đủ, Aqua chia hai số vừa tính:``\textit{Vậy xác suất là \(\frac{3}{10}\)!}'' Cô hí hửng dò đáp án,
nhưng rồi ngớ người ra khi nhìn cả bốn đáp án đó trong đề\dots

\begin{table}[H]
    \begin{tabular}{>{\arraybackslash}p{0.25\linewidth}>{\arraybackslash}p{0.25\linewidth}>{\arraybackslash}p{0.25\linewidth}>{\arraybackslash}p{0.25\linewidth}}
        A. \(\frac{1}{6}\) & B. \(\frac{1}{5}\) & C. \(\frac{2}{5}\) & D. \(\frac{1}{3}\) \\
    \end{tabular}
    \caption{Caption}
\end{table}

\dots cô ngớ người ra, trán bắt đầu lấm tấm mồ hôi: ``\textit{Gì chứ!? Không có đáp án nào là \(\frac{3}{10}\) hết! Chuyện gì thế này!?}''

Cô quay lại kiểm tra, nhưng càng làm, kết quả vẫn là \(\frac{3}{10}\). Aqua bắt đầu hoảng loạn, tay vò
tóc: ``\textit{Không lẽ mình làm sai đâu đó? Nhưng mình đã tính rất cẩn thận rồi mà!}''

Kuon đang đi tuần, chú ý đến nàng Hầu gái đang hoảng, khẽ nhướng mày: ``Làm sao thế,
Aqua-chan? Trông em căng thẳng nhỉ?''

Aqua hướng về cậu con trai, tỏ vẻ ngỡ ngàng với tình cảnh éo le hiện tại của mình: ``Kuon-senpai! Em\dots\ em tính mãi mà không ra được đáp án trong bài!''

Kuon nhìn thoáng qua bài làm của cô, lắc đầu cười khẽ: ``Em đúng rồi mà, chỉ là\dots\ anh
cố tình không để \(\frac{3}{10}\) trong đáp án thôi.''

Nàng Hầu gái trố mắt ra: ``C-Cái gì cơ!?'' Cậu Tổng phì cười, giải thích rõ hơn: ``Đáp án
của em chính xác. Nhưng em có lẽ không đọc kỹ đề bài rồi. Đề người ta hỏi là lấy giá trị gần nhất mà, có phải giá trị chính xác đâu. Con số mà em ra gần nhất với kết quả \(\frac{1}{3}\) đó.''

Aqua như bị sét đánh ngang tai, gương mặt đỏ bừng: ``R-Rút gọn!? Sao trong đề không
ghi rõ hơn chứ!''

Cậu khoanh tay, điềm tĩnh trả lời: ``Đề bài ghi rõ mà, chắc là căng thẳng quá nên không tìm ra. Em nên nhớ, bài kiểm tra cũng là để kiểm tra sự cẩn thận và tư duy logic của mình nữa.''

Nàng Hầu gái run tay, chỉnh lại bài làm của mình với phản ứng vừa uất ức vừa sợ hãi:
``Rút gọn\dots\ trời ơi, bài toán này thật sự là ác mộng mà\dots''

A-chan nhìn cảnh tượng Akutan đang lẩy bẩy, khẽ thở dài: ``Lại thêm một người bị
Hikawa-san đánh đòn tâm lý rồi\dots'' Choco bật cười, đùa nhẹ: ``Không biết đây là kỳ
thi hay buổi tra tấn tinh thần nữa\dots''

\ct

Ở một diễn biến khác, Matsuri đang làm bài thi với vẻ mặt căng thẳng, phần vì bài toán
quá khó, phần vì đang phải chịu một cơn mót tiểu đang mạnh mẽ ập đến.

Cô cố gắng tập trung vào bài toán, nhưng nhìn xuống câu hỏi tiếp theo trên tờ đề, sự không
thoải mái lại càng khiến cô không thể tập trung vào những con số, thậm chí còn khiến cô
hoang mang hơn, khi nó không khác gì đổ thêm dầu vào lửa.

Không gì khác, đó là câu hỏi yêu cầu tính thể tích của một bồn nước với hình ngoằn ngoèo
được cho sẵn, sử dụng phép tích phân ba lớp. Đề bài như sau: ``Một bồn nước có hình
dáng như một chiếc phễu, chiều cao  \(h\) và bán kính đáy \(r\), và các chỉ số cho hình bên. Tìm công thức đúng để tính thể tích của bồn nước bằng cách áp dụng phép tích phân ba lớp (ví dụ: \(\iiint r^2\sin(\theta)drd\theta d\phi\).)''

Nàng Cổ động nhìn vào bài toán với ánh mắt mơ màng, một tay run rẩy cầm bút, một tay
giữ lấy hạ bộ của mình. ``\textit{Cái này là gì nhỉ\dots\ Tích phân ba lớp là sao? Mình không nhớ
    học cái này bao giờ\dots}''

Từ phía xa, Kuon đang dõi theo cô nàng đang tuyệt vọng làm toán trong khi đang chiến
đấu với ``cơn lũ'' có thể ``ập'' tới bất cứ lúc nào. Cậu bật cười: ``Mót vệ sinh thì ít nhất
cũng phải lên tiếng chứ.''

``Có vẻ như Matsuiro-san không để ý tới lưu ý mà cậu nói nhỉ,'' A-chan thủ thỉ vào tai cậu,
tỏ vẻ lo ngại cho cô nàng.

``Chắc thế rồi,'' cậu Tổng chép miệng.

Với cái bụng đang ``kêu gào'', cô vội vàng cố làm một phép tính đại khái. Cô lấy công
thức thể tích hình trụ là \(V = \pi r^2h\), nhưng do không chắc chắn, cộng với tinh thần đang
hoảng loạn, cô tự nhẩm: ``\textit{Chắc là\dots\ như thế này thôi\dots\ Tích phân ba lớp thì sao, tính ra
    hình trụ chắc là xong rồi\dots}''

Cô nàng nhìn nhanh qua đáp án, nhưng với sự tiến thoái lưỡng nan hiện tại, cô liền nhắm
mắt khoanh bừa đáp án D.

Ngay lúc đó, ``cuộc gọi thiên nhiên'' của cô ngày càng lúc mạnh mẽ hơn, khiến cô không
thể ngồi im thêm được nữa. Cô liền đứng dậy, tiến về phía cửa phòng thi mà không hề
nhìn thấy giám thị chính đang đứng gần đó.''

Cậu Tổng thấy vậy, lên tiếng hỏi: ``Đi đâu vậy?'' Matsuri, mặt đỏ như gấc, trả lời nhanh
chóng: ``E-Em đi vệ sinh!''

Cậu nhìn nàng Lễ hội với ánh mắt hơi ngạc nhiên, nhưng rồi sau đó đáp lại: ``Ít nhất thì
cũng phải lên tiếng chứ.'' Cô nàng vội xin lỗi, rồi chạy thẳng ra ngoài, thở phào nhẹ nhõm
khi được ``giải quyết'' kịp thời.

Cậu tiến đến chỗ ngồi của cô nàng, nhìn vào tờ bài vẫn đang làm dở. Cậu thở dài: ``Đi
mà không che lại bài kìa. Chán thật chứ.''

Cậu liếc về câu hỏi mà cô vừa chọn bừa đó, hơi trố mắt ngạc nhiên. ``\textit{Quái lạ, em ấy chọn
    đúng rồi kìa. Đây là một trong những câu khó của đề đấy.}''

Cậu khẽ mỉm cười, rồi bước đi kiểm tra các thí sinh khác. Đúng lúc đó, Matsuri trở lại
phòng thi, nhìn vào bài làm của mình. Cô nàng khẽ thở dài, nhìn lại câu hỏi tìm công thức
của mình: ``\textit{Chắc là mình làm sai rồi. Thôi kệ.}''

Cô khẽ chẹp miệng, chấp nhận mất điểm câu hỏi đó, rồi chuyển sang câu hỏi tiếp theo.

Từ phía xa, Kuon nhìn cô nàng đang tập trung làm các câu khác, khẽ mỉm cười: ``\hl {Có lẽ
    em ấy sẽ vui lắm khi biết em ấy vừa trả lời đúng đấy.}''

A-chan nghe vậy, bật cười: ``\hl{Cô ấy lại khoanh đại mà lại đúng rồi, thật đúng là vận may\dots}''

Choco ở cạnh trêu chọc thêm: ``Có khi nào Kuon-sama cố tình tạo ra những bài tập khó
để kiểm tra tâm lý các em không nhỉ? Để họ khoanh đại cho vui?''

Kuon nhún vai, không đáp lời hai người họ, rồi tiếp tục đi tuần.

\ct

Sau khi đã làm bài gần hết thời gian, Roboco bất ngờ úp bài thi lại rồi ngả người lên bàn,
nhắm mắt. Hành động này khiến A-chan nhìn cô với vẻ khó hiểu.

``Roboco-san, cô làm xong rồi à?'' nàng Đeo kính hỏi trong sự ngạc nhiên.

Roboco mở mắt một chút, nhìn vào tờ bài thi và trả lời một cách dửng dưng: ``Em xong
rồi ạ\dots''

Nàng Đeo kính nhíu mày với sự bất cẩn, vẫn cảm thấy chưa yên tâm: ``Tốt nhất là cô nên
kiểm tra lại bài đi.''

Nàng Người máy lắc đậu, vẻ mặt quyết đoán: ``Thôi không cần. Em nghĩ mình đã làm
đúng rồi.''

Kuon ngay ở bên cạnh A-chan, thở dài: ``Đừng cậy mình là người máy mà có thể chủ
quan, PON ạ. Máy móc nhiều khi còn có thể tính sai đấy.''

Mặc cho những lời cảnh cáo, Roboco vẫn nhún vai, tỏ vẻ rất tự tin: ``Chắc gì em bị như
thế\dots''

Nghe câu trả lời có vẻ như ``mặc kệ sự đời'' của cô A-chan và Kuon chỉ thở dài, không nói
gì thêm. Họ biết rằng việc này giống như một cuộc trò chuyện không có điểm dừng với Robo-PON - cô máy tính thông minh, nhưng đôi lúc lại có chút ngớ ngẩn. Nàng ``Tăng
ca'' nhìn Roboco nằm ngủ một cách ung dung, không hề lo lắng gì về bài thi nữa.

Choco ở phía xa, chỉ biết lắc đầu thầm nghĩ: ``\textit{Cái cách mà cô ấy làm bài thi\dots\ rõ ràng là
    một thách thức đấy. Hy vọng rằng cô ấy sẽ không quá ngạc nhiên với kết quả của mình\dots}''

Giữa không gian yên tĩnh của phòng thi, A-chan, Choco và Kuon tiếp tục đi tuần qua lại,
mắt vẫn theo dõi các thí sinh làm bài, nhưng thỉnh thoảng lại lén nhìn Roboco đang ngủ
say. Cô nàng người máy ấy chẳng hề biết, trong thế giới thi cử, có những thứ không phải
lúc nào cũng phụ thuộc vào độ chính xác của phần cứng, mà kiểm thử mới là thước đo
của sự thành công.

\ct

Giờ làm bài đã gần hết, không khí trong phòng thi trở nên căng thẳng hơn bao giờ hết.
Những tiếng bút viết lách đều đặn vang lên, nhưng càng về cuối, tiếng bút chậm dần lại.
Cả căn phòng như chìm trong sự im lặng, chỉ còn lại tiếng đồng hồ tích tắc từng giây,
nhắc nhở các thí sinh về thời gian còn lại.

A-chan và Choco cũng bắt đầu thu dọn các giấy tờ, mắt nhìn đồng hồ rồi nhìn lại các thí
sinh một lần nữa. Lúc này, Kuon đứng dậy, vỗ tay gây chú ý: ``Được rồi, hết giờ làm bài,
mọi người kẹp đề và bài làm vào đề thi rồi chuyển ra đầu bàn nào.''

Nghe lệnh, các cô gái ngay lập tức làm theo, cẩn thận xếp lại bài thi của mình. Vừa chuyển
bài thi, họ vừa nói chuyện với nhau về bài thi họ mới làm. Lúc này, cả ba giám thị bắt
đầu thu bài của các thí sinh thi.

Nhưng không phải ai cũng vội vã hoàn thành như thế. Từ góc phòng, Ayame vẫn đang
ngồi nghiêng đầu, tập trung viết, không hề có dấu hiệu dừng lại. Cứ mỗi khi cậu nhìn về
phía cô, thấy nàng vẫn lẩm bẩm tính toán trong đầu, cậu Tổng lại không khỏi thắc mắc.

Cuối cùng, khi không thể để tình huống tiếp tục như vậy, cậu quản lý liền đi đến gần nàng
Quỷ sừng, nhẹ nhàng nhắc nhở: ``Ojou-san, hết giờ rồi. Nộp bài thôi.''

Nàng Quỷ sừng giật mình, nhìn lên cậu Tổng đang sừng sững đứng đó, rồi lại nhìn xuống
bài thi của mình, bối rối đáp: ``Ơ từ từ, em còn mấy câu chưa làm nữa\dots''

Kuon mỉm cười, nhưng ánh mắt hơi sắc bén hơn thường lệ: ``Còn mấy câu nữa?''

Nàng Oni hít một hơi thật sâu, rồi thở dài, lí nhí nói: ``Mười câu nữa\dots''

Tới đây, cậu không thể giữ được vẻ mặt bình thản nữa, chỉ có thể đảo mắt nói nhanh:
``\dots Hoặc là nộp bài, hoặc là cho trứng ngỗng. Nhanh lên.''

Nàng Quỷ sừng không thể làm gì khác ngoài việc nén lại sự thất vọng, hậm hực xếp bài
lại và đưa cho Kuon, mặt đầy vẻ bất lực. Cậu con trai cười nhẹ, thu hết tất cả bài thi, rồi
bắt đầu xách chúng đi về phía bàn giáo viên.

Sau khi Kuon vừa đi về phía bàn giáo viên, không khí trong phòng bắt đầu rộn ràng hẳn
lên. Dù cậu đã nhắc mọi người giữ trật tự, mười đám cô gái vẫn không thể không trao đổi
với nhau về cảm nhận của mình sau bài thi đầy ``thử thách''.

Mỗi người lại có một cảm xúc khác nhau, nhưng đại đa sổ là mệt mỏi, bất bình, thậm chí
nổi cáu với bài thi đầy hóc búa này.

``Đề lần này thật sự khó hơn mình nghĩ\dots\ Có vài câu mình không chắc lắm,'' Sora vừa nói
vừa khẽ cười, cố gắng động viên mọi người xung quanh.

Roboco vừa mới ngủ dậy, cô nàng duỗi cơ, chẳng mấy quan tâm đến việc mọi người đang
thảo luận sôi nổi: ``Ể? Xong hết rồi à? Mình chẳng nhớ gì ngoài việc ngủ ngon cả.''

Miko xoa đầu, trông khá bối rối, rồi cố gắng cười trừ để che đi sự lúng túng của mình:
``Tôi á, ne? À thì, hầu hết các câu trong đề mình chọn đều theo\dots\ thần linh (Kami-sama)
mách bảo thôi, ne!''

Fubuki thì vò đầu bứt tai, rõ ràng chưa hết ám ảnh bởi các phép toán khó nhằn, mặc cho
cô nàng đã tỏ vẻ tự tin trước đó: ``Trời ơi, tính toán đầu óc muốn nổ tung luôn! Cái gì mà
tích phân, mà đạo hàm\dots\ Kuon-senpai ác quá đi!''

Haato hăng hái giơ tay lên, nhưng nụ cười của cô hơi gượng, như thể cô cũng vừa trải qua
một thử thách cam go: ``Thú vị nhỉ!? Đề này giống như\dots\ một trận chiến sống còn với số
học ấy!''

Matsuri thì lắc đầu cười trừ, nhưng ánh mắt trông như vẫn đang loay hoay với số phận
của mình: ``Hồi nãy mình còn khoanh lụi vì\dots\ mót tiểu. Chắc là câu đó cũng không quan
trọng đâu nhỉ? Cùng lắm mất 0.2 điểm thôi.''

Mel mỉm cười dịu dàng, cố gắng giữ thái độ tích cực dù bên trong có chút hoang mang:
``Một vài câu mình không biết cách làm, nhưng đề này cũng khá thú vị đó chứ.''

Aki cầm cốc nước trên tay, khẽ thở dài như thể còn đang mơ về việc nhảy múa giữa cánh
đồng: ``Mình chỉ muốn biết tại sao đề lại dài như vậy\dots\ Mình chưa từng làm bài thi nào
nhiều câu đến thế.''

Shion bực bội, vắt óc suy nghĩ lại mấy câu khó nhằn, dù giờ chẳng còn ý nghĩa nữa:
``Hừm, sao có câu làm đi làm lại mãi mà chẳng ra đáp án vậy!? Đề này chắc chắn có bẫy,
đúng không!?''

Ayame thì tỏ vẻ hậm hực hơn, rõ ràng vẫn chưa nguôi giận vì phải nộp bài trong khi còn
tới mười câu chưa làm xong: ``Không công bằng! Mình còn chưa làm hết nữa\dots\ Sao
không cho thêm thời gian chứ, yo!?''

Aqua ôm đầu, ánh mắt như vừa bị phản bội bởi chính chiếc máy tính và những con số:
``Mình ghét nhất là mấy câu mà đáp án mình tính ra không trùng với bất kỳ lựa chọn nào!''

Subaru thì xoa cằm, ra vẻ ngẫm nghĩ, nhưng trông có vẻ đang cầu nguyện hơn là phân
tích: ``Đề này kiểu như muốn mình điên luôn ấy! Câu nào mình cũng thấy nghi nghi,
nhưng chọn xong rồi thì\dots\ thôi vậy.''

Mio gục mặt xuống bàn, vừa thở dài vừa than thở không ngừng: ``Kuon-senpai\dots\ Anh có
chắc đây là đề dành cho người thường không? Cả buổi đầu em chỉ lo tính giải tích với
tích phân thôi á\dots''

Korone gãi đầu, trông như đang đấu tranh giữa sự mệt mỏi và quyết tâm không chịu thua:
``Công nhận đề dài thật. Tôi làm gần hết nhưng vẫn còn vài câu\dots\ toàn mấy từ chuyên
ngành nghe lạ hoắc á, Ogayu!''

Okayu khoanh tay, cười nhẹ nhưng không giấu được chút bối rối: ``Bình tĩnh nào, Korosan.
Tôi cũng chọn đáp án dựa vào giác quan mèo, có gì sai thì đổ tại cảm nhận vậy.''

Rushia lẩm bẩm, cúi gằm mặt, tay nghịch nghịch cái máy tính của mình: ``Mấy câu mình
làm đúng có lẽ là nhờ may mắn thôi. Chứ kiến thức như thế này thì làm sao mình biết
được!''

Noel nhăn mặt, nhìn về phía bài làm ở trên bàn giáo viên của mình, rõ ràng vẫn còn chưa
hài lòng: ``Có phải Kuon-senpai đang cố ý thử thách trí thông minh của bọn mình không?
Nếu đúng thế thì hơi quá đà rồi đó.''

Pekora thở hổn hển, vừa nói vừa xoa đầu mình như thể đang cố làm nguôi sự hoảng loạn:
``Đề này khó thế này thì ai mà làm được chứ, peko!? Nếu mình qua được bài này thì chắc
chắn là nhờ may mắn thôi!''

Riêng chỉ có cô em gái Ryuuhou của cậu là tỏ vẻ rất hài lòng với bài thi của mình: ``Mình
dám chắc là mình có thể qua được điểm trung bình! Cơ mà\dots'' Cô tỏ vẻ lo ngại với bài
làm của mình, không hoàn toàn tự tin với bài làm của mình, rồi quay sang nhìn phản ứng
mãnh liệt của mọi người, cười khổ: ``Dám chắc là đề thi của họ trông giống như đề thi
học thuật cho bên Bách Khoa vậy á\dots\ Onii-san đúng là chơi khó thật.''

Tất cả những lời bàn tán này rộn ràng khắp phòng, khiến Kuon chỉ có thể cười trừ, tỏ vẻ
đồng cảm với họ. ``\textit{Mình phải thừa nhận là đề thi của mình có phần hơi quá sức với họ
    thật\dots\ Chẳng biết họ đã chịu đựng bao nhiêu nỗi đau như thế nữa.}''

Cậu tập hợp tất cả bài lại rồi ngồi bàn giáo viên, vỗ tay gây chú ý tới mọi người, chuẩn bị
chấm bài.

A-chan nhìn xấp bài mà không có bút, thấy làm lại hỏi: ``Cậu có mang bút gì để chấm
không?''

Kuon cười, lắc đầu: ``Không cần chấm tay đâu. Em có cách này.'' Nói rồi, cậu lấy điện
thoại ra và mở một ứng dụng chấm điểm.

Cậu đặt trước phiếu trả lời trắc nghiệm của Sora, để cho máy tự quét đáp án. Chỉ mất
chưa đến ba giây, ứng dụng đã quét xong mã đề và các đáp án cô chọn. Màn hình hiện
lên một con số.

Cậu Tổng đọc to số điểm được viết trên đó: ``Tokino Sora, 8.6 (tám chấm sáu) điểm!'' Fubuki ngạc nhiên
với tốc độ chấm điểm của cậu: ``Nhanh thế!''

Nàng Đeo kính nhìn số điểm được ghi trên máy cậu, tỏ vẻ hơi nghi ngại: ``Cậu có chắc
đây là điểm đúng không?''

Kuon đáp với sự tự tin: ``Ứng dụng này là do em lập trình, nên chắc chắn là không sai
được đâu.''

Trong khi đó, Sora vừa nghe điểm số của mình, khẽ tỏ vẻ ngạc nhiên: ``8.6 ạ? Sao em
cảm thấy có gì đó hơi sai\dots''

Nhìn thấy phản ứng khó hiểu của cô ấy, cậu giải thích thêm: ``Hình như A-san nói rằng
bọn anh chấm theo thang điểm 10 đúng không nhỉ?''

Nàng Thần tượng ngẫm lại những gì mà cô nàng Đeo kính đã nói trước khi vào phòng thi,
cô gật đầu. ``Vâng. Thì\dots''

``Nghĩa là em đạt được 86/100 điểm đấy. Chúc mừng em.''

Sau khi biết được điểm chính xác của mình, cô nàng tỏ vẻ mừng rỡ, đôi mắt sáng lên vì
đã qua được bài thi một cách ngoạn mục: ``Yatta! Vậy là tốt rồi!''

``Đồng nghĩa với việc là xếp loại 1. A-san, chị ghi điểm này vào tờ danh sách giúp em
nhé.''

Cậu đưa danh sách điểm các thí sinh đã tham gia cuộc thi này cho người đồng nghiệp.
A-chan gật đầu, vừa viết vừa lẩm bẩm: ``Nhìn nhanh gọn thế này thì tiện lợi thật.''

Choco đứng cạnh cô, cũng mỉm cười gật đầu đồng tình.

Cứ thế, Kuon tiếp tục quét từng bài làm, tạo nên sự hào hứng lẫn hồi hộp trong căn phòng.

\ct

Tất nhiên, cũng không tránh khỏi những trường hợp dở khóc dở cười, kiểu như\dots

Khi cậu Tổng chấm bài của Noel, thấy có những câu trả lời vẫn bỏ trắng, cậu ngạc nhiên
hỏi: ``Noel-chan, sao em lại có nhiều câu chưa tô thế này?''

Nàng Hiệp sĩ nghe vậy, chớp mắt tỏ vẻ vô tội: ``Anh bảo là câu nào không làm được thì
cho qua mà.''

Lời giải thích của cô khiến cho cậu thở dài ngán ngẩm, gõ nhẹ vào bàn: ``Cho qua là để
tí nữa quay lại làm ấy. Còn nếu không biết thì cứ khoanh bừa một đáp án đi, biết đâu lại
đúng.''

Subaru cũng gật đầu phụ họa: ``Chị cũng làm như thế đấy, phần lớn câu của anh ấy ra khó
quá mà.''

Mio vội cắt ngang, cẩn trọng nói: ``Nhưng cái đó may rủi lắm. Chỉ dùng nó như cứu cánh
cuối cùng thôi.''

Những lời giải thích từ nàng Sói đen, nàng Vịt và cả cậu Tổng khiến Noel trố mắt ra ngỡ
ngàng: ``Thế à? Em không biết đấy\dots''

Kuon vỗ trán, ngán ngẩm lắc đầu với suy nghĩ giản đơn tới khờ dại của cô nàng. Cậu cầm bài của cô
ấy, quét nhanh và đọc to số điểm một cách dứt khoát: ``Chính vì sự PON của em mà em
chỉ được 2.8 điểm, ứng với 28/100 thôi. Loại 4.''

Nghe tới đây, nàng Hiệp sĩ đờ người, rồi ôm đầu kêu lên đầy đau đớn, khi cô đã ở rất, rất
sát với loại 3: ``Tại sao chứ!? Em chỉ cần thêm một câu nữa là thoát cửa tử mà\dots''

Biểu cảm của cô nàng khiến A-chan mỉm cười khó xử, cố gắng nhẹ giọng xoa dịu: ``Cũng
không hẳn là thế đâu, có thể lần sau sẽ\dots\ cải thiện được.''

Tuy nhiên, trước khi kịp dứt lời, Pekora chen vào với nụ cười ranh mãnh: ``Bà đúng là gà
quá, peko! Đồ ngu! Ngờ-u-ngu!''

Kuon ngước lên, nhìn cô nàng Thỏ tinh quái kia đang giở giọng sỉ vả, tỏ vẻ không đồng
tình: ``Em cũng chỉ có được 64 điểm thôi đấy. Đừng có đứng đó mà khịa người ta.''

Nàng Thỏ đáp lại bằng một tràng cười đầy hả hê: ``Còn hơn là chịu phạt do bà ấy tự đặt
ra á, peko! Em lại còn được tiền nữa chứ, ha \uparrow ha \downarrow!''

Nghe những tràng cười đầy xấu xa và mỉa mai đó, nàng Hiệp sĩ mím môi lại, giọng run
run: ``Bà đúng là quá đáng\dots\ Tôi hận bà, Pekora!''

Choco, người đang quan sát từ phía sau, cười khúc khích rồi xen vào: ``Pekora-sama, sỉ
nhục người khác không hay ho gì đâu. Noel-sama à, không cần cãi nhau làm gì. Ai biết đâu được, biết đâu lần sau am ấy sẽ giỏi hơn, nhờ biết cách `khoanh bừa' đúng quy tắc?''

Cậu Tổng gật gù đồng tình: ``Đúng đấy. Lần sau thông minh hơn một chút đi, đừng có để
óc não nho nữa. Đấy là anh còn chưa hỏi em chuyện tại sao lại có hình phạt như sỉ nhục
người khác do em với Shion-chan bày ra đâu.''

Noel chỉ biết cúi đầu, thở dài trong sự tiếc nuối, còn Pekora thì vẫn ôm bụng cười khanh
khách.

\ct

Một trường hợp khác, khi cậu Tổng đang chấm bài cho Roboco. Ngay khi cậu nhận được
điểm của cô nàng, cậu trố mắt ra ngạc nhiên, đến mức cậu phải chấm đi chấm lại.

Nhưng rồi, cái kết quả ấy vẫn xuất hiện, không lệch một điểm nào. A-chan và Choco nhìn
vẻ mặt ngỡ ngàng của cậu, tò mò hỏi:

\begin{dialogue}
    \speak{A-chan} Sao thế, Hikawa-san?
    \speak{Choco} Có chuyện gì mà anh trông ngạc nhiên vậy?
\end{dialogue}

Kuon không nói gì, liền giơ điện thoại của cậu ra để xem số điểm mà Roboco nhận được.
Nàng Y tá thì vỗ trán, lắc đầu thất vọng, còn A-chan thì thở dài, nói: ``Tôi đã bảo rồi mà.''

Kuon nhíu mày, gọi lớn cô nàng Người máy đang trò chuyện với Okayu và Korone: ``Đại
PON ơi, anh thấy ngạc nhiên với số điểm của em đấy.''

Giống như một thói quen, Roboco ngẩng mặt lên, thốt lên: ``Đã bảo rồi, em không có
PON! Em là rô-bốt công nghệ cao!''

Cậu nhìn lại bài của cô nàng, khẽ bật cười với một trong những kết quả mà cô làm. Cậu
nói một cách chậm rãi, nhấn nhá từng chút một: ``Vâng, công nghệ cao đến mức mà \u=line{chẳng có câu nào} trong này làm đúng cả. Thế là thế nào đây?''

Roboco giật thót mình, há hốc miệng với kết quả đầy ngỡ ngàng này: ``\dots Ơ? Rõ ràng là
em tính đúng rồi mà.''

``Đúng cái nỗi gì. Sai be sai bét hết cả lên này.'' Cậu con trai đứng dậy, giơ tờ giấy trước
mặt cả lớp, giọng cười như không thể tin nổi: ``Cả lớp nghe anh đọc ngay ở câu hỏi đầu
tiên nhé: Thực hiện phép tính \(3 + 4\). Và em ấy chọn kết quả bằng 8 ạ!''

Cả phòng học cười ầm lên, tỏ vẻ sốc đến bật ngửa trước câu trả lời sai đến mức ngớ ngẩn
đến thế. Chủ nhân của câu trả lời đó đỏ mặt, đập bàn cố gắng minh oan: ``Em tưởng cái
đó là đúng mà! Em đã tính bằng siêu máy tính trên não của đó!''

Cậu Tổng nhướng mày, cố gắng tỏ vẻ nghiêm túc nhưng không thể: ``Đúng cái nỗi gì. Về
xem lại cái não của em đi. Với cả cái máy tính mà em mang vào chỉ để làm cảnh thôi à?''

A-chan nhìn cô nàng đang rơm rớm nước mắt vì kết quả trái ngược với mong muốn, hoàn
toàn đối lập với vẻ tự tin ban nãy, cũng chỉ có thể thở dài bất lực: ``Tôi đã bảo rồi mà
không nghe. Máy móc chưa chắc chính xác đâu mà.''

Roboco méo mặt, dường như không thể chấp nhận với kết quả này: ``Tại sao chứ!! Tại
sao!?'' Mặc cho cô rầy la cả phòng, cậu Tổng vẫn đọc điểm và ghi phán xét một cách tàn nhẫn:
``Robo-PON, một quả trứng ngỗng, xếp loại 4. Và đừng cãi, cái tên này hợp với em thật
đấy.''

Nàng Người máy bực bội, dậm chân xuống sàn ăn vạ: ``Em không có PON thật mà!!'',
trong khi mọi người xung quanh đều đang cười trước sự đau khổ của cô.

\ct

Tới bài của cặp đôi Okayu-Korone. Ngay khi cậu vừa chấm xong bài của họ, Kuon ngước
nhìn lên, gọi cặp đôi đang trò chuyện lên.

\begin{dialogue}
    \speak{Okayu} Chuyện gì thế, Kuon-senpai?
    \speak{Korone} Điểm của bọn em bất thường quá hở?
\end{dialogue}

Cậu nhìn họ với vẻ mặt không mấy hài lòng, rồi nghiêm nghị nói: ``Anh để ý suốt thời
gian làm bài, thấy Okayu-chan và Korone-chan có những hành động mờ ám lắm đấy.''

Nàng Mèo nghe vậy, giả vờ không hiểu cậu đang nói gì, đáp lại: ``Em không hiểu anh
đang nói gì cả.''

A-chan, người cũng đã để ý tới cặp đôi này trong suốt thời gian kiểm tra, đưa cho họ
những mảnh giấy nhỏ ghi đáp án trên, thở dài: ``Tôi có nhìn thấy vài mảnh giấy trong
ngăn bàn của hai người. Hai người trao đổi bài trong giờ kiểm tra à?''

Cả Okayu và Korone khẽ giật nảy mình, rồi xoa gáy cười khì tỏ vẻ vô tội.

\begin{dialogue}
    \speak{Korone} Bại lộ rồi, hihi\dots
    \speak{Okayu} A-chan đúng là mắt tinh thật á\~{}
\end{dialogue}

Kuon thở dài, gằn giọng với thái độ lồi lõm của cặp đôi: ``Hì cái gì mà hì. Mỗi đứa trừ
25\% điểm bài thi. Vi phạm quy chế thi rồi đấy!''

Nghe phán quyết như vậy, cặp đôi bắt đầu giương ra nét mặt đáng thương, với đôi mắt
long lanh và vẻ mặt đượm buồn, giọng nải nỉ:

\begin{dialogue}
    \speak{Okayu} Kuon-senpai! Tha cho bọn em lần này đi mừ\~{}
    \speak{Korone} Đúng á! Bọn em chi nhìn nhau một xíu thôi, không tính là quay cóp đâu!
    \speak{Okayu} Bọn em hứa là lần sau sẽ không như thế nữa đâu, meo\~{}
\end{dialogue}

Nhìn vẻ mặt giả bộ đáng thương của hai người họ, Kuon khẽ bật cười. Dù thế, cậu vẫn
nghiêm nghị nóiđầy kiên quyết: ``Đừng có nịnh. Vi phạm vẫn là vi phạm. Anh không
phạt nặng là may cho hai đứa rồi đó.''

A-chan ở ngay bên cạnh, thở dài: ``Tôi đã phổ biến quy chế ngay từ đầu rồi mà hai cô có
chịu nghe đâu. Thế vẫn còn hơn là nhận con 0 đấy.''

Choco thì nhìn họ với vẻ thất vọng: ``Gian lận trong thi cử không phải là điều tốt đâu,
Korone-sama, Okayu-sama. Mau xin lỗi anh ấy đi.''

Cặp đôi thở dài, cúi đầu xin lỗi. Nàng Miêu quay sang trách yêu bạn mình: ``Tại bà không
giấu mấy mảnh giấy cẩn thận đấy.''
Nàng Khuyển lè lưỡi, gật đầu: ``Thì, tôi đã xin lỗi rồi mà, gâu gâu!''

\ct

Đỉnh điểm là Miko ``ưu tú'' và ``bà hoàng tính toán'' Fubuki, khi họ thậm chí còn sai ở
những câu hỏi rất cơ bản.

Cậu Tổng cầm phiếu trả lời của Miko, gật gù khen ngợi: ``Anh đánh giá cao khi những
câu siêu khó như định thức ma trận, tổ hợp hay xác suất thống kê thì em làm đúng đấy.''

Nàng Vu nữ nghe những lời có cánh từ cậu, liền ngẩng cao đầu, cười toe toét: ``Thấy
chưa! Miko đúng là ưu tú mà lị! Yatta\~{}!''

Nhưng rồi, cậu nhíu mày, đổi hẳn sang giọng nghiêm khắc: ``Thế tại sao mấy câu cơ bản
như cộng trừ nhân chia lại sai lên sai xuống thế?''

Lời nói của cậu Tổng như một gáo nước lạnh vào niềm kiêu hãnh của cô, làm cho cô sững
người lại, chớp mắt liên tục: ``\dots Ể? Anh đang nói gì chứ, ne?''

Cậu đặt bài xuống, cầm phấn lên bảng viết một phép tính mà cô làm sau trên đề và nói:
``Miko-chan, em thử làm lại phép tính \(5 + 36 \div 3\) xem nào.''

Cô nàng hớn hở bước lên bảng, tay thoăn thoắt viết lời giải: ``Thì chẳng phải là phép toán
này làm từ trái sang phải sao? 5 cộng với 36 là 41, rồi chia cho 3 thì được 13 à?''

Kuon gật đầu, khoanh tay nhìn cô nàng đang hớn hở làm. ``Tiếp đi.'' Tất nhiên, ngay lập
tức đã có những tiếng xì xào: ``Sao tôi cảm thấy có gì đó sai sai nhỉ\dots''

Mặc dù vậy, Miko vẫn tiếp tục trình bày bài làm của mình, không hề nao núng: ``Còn đuôi
,6667 thì Miko nghĩ nó là dư 2 vì chỉ có đáp án này thôi ne. Đúng chưa?''

Những tiếng xì xào ngày càng rộ lên, thậm chí còn dáng dấp những tiếng nhắc từ bên
dưới: ``Sai rồi kìa Miko-chan!''

Cậu nhìn vẻ mặt tự đắc của nàng Vu nữ, phì cười nhưng cũng thương cảm: ``Đúng là như
thế, nếu theo đáp án em chọn.'' Để rồi, cậu lắc đầu ngán ngẩm: ``Nhưng mà kết quả em
ra sai rồi đấy. Trừ thêm một điểm vì lập luận sai kiến thức cơ bản nha.''

Miko há hốc miệng, quay sang nhìn cậu với vẻ ngạc nhiên: ``Sao lại thế, ne!? Rõ ràng
chẳng phải phép tính phải thực hiện từ trái sang sao!?''

Từ bên dưới, Subaru cười trừ với độ ngốc nghếch của tiền bối mình: ``Miko-chan ơi, phép
tính đó phải thực hiện nhân chia trước, cộng trừ sau mà\dots''

Lời giải thích của nàng Vịt khiến cô nàng im bặt, ánh mắt vẻ đờ đãn như vừa hiểu ra điều
gì. ``\dots Ồ, ra là vậy, ne.''

Ryuuhou chêm vào, tỏ vẻ điềm tĩnh: ``Kiến thức này là kiến thức cấp tiểu học luôn đó,
Miko-onee-chan ạ\dots\ Em tưởng cái này chị phải nhớ chứ?''

Cậu Tổng cười nhạt: ``Em trông mong gì từ bộ não của một em bé chứ, Ryuuhou? Em ấy thậm chí còn chưa tốt nghiệp mẫu giáo mà.''

Vừa dứt câu, nàng Vu nữ phản pháo lại: ``A-Anh đang nói gì!? Miko không phải là em
bé! Miko ưu tú mà!''

Bỏ mặc hoàn toàn những lời ca thán của cô, cậu quay sang Fubuki, giơ bài làm của cô
lên: ``Fubu-chan cũng mắc sai lầm y chang như vậy. Lúc anh đi tuần anh thấy em thực
hiện phép tính cũng sai như thế.''

Nàng Cáo trắng cười gượng, vẫy tay như tự biện hộ: ``E-Em vốn dĩ ghét môn toán mà\dots
Mấy cái này không hợp với em chút nào\dots''

Kuon thở dài, bật cười nhẹ: ``Nhưng mà được cái là em còn biết dùng máy tính kiểm tra
lại kết quả. Còn Ê-lít này thì tự tin với kết quả của mình quá.''

Cả lớp cười ầm lên trước câu đánh giá đầy ẩn ý của cậu, đồng thời chọc ghẹo cô nàng Vu
nữ ấy. Cô nàng phồng má lên, giậm chân ăn vạ như em bé: ``Miko không thích bị so với
người khác đâu! Đừng có chọc tôi chứ, ne!!''

Như để xát muối vào nỗi đau của cô, Choco bất ngờ lên tiếng, nở một nụ cười dịu dàng:
``Vậy lần sau nhớ ưu tú ở mọi phép tính luôn nhé, Miko-sama\~{}''

Nàng Vu nữ vùng vằng quay đi, lí nhí nói trong sự cay cú: ``Miko sẽ làm tốt hơn\dots\ Lần
sau nhất định sẽ ưu tú hơn!''

\ct

Thậm chí đến cả người nhà của Kuon cũng không thể thoát khỏi lưỡi hái tử thần của cậu.
Cô em gái của cậu cũng bị Kuon chỉ ra lỗi sai ngay trước cả lớp.

Cậu Tổng nhìn vào bài làm của cô em gái, rồi chậm rãi hỏi: ``Ryuuhou, em có nghĩ mình
làm sai ở đâu không?''

Ryuuhou ngẩng mặt lên, nói với anh trai đầy tự tin: ``Không ạ. Em nghĩ mình làm khá tốt
đấy chứ!''

Cậu nhìn lại bài làm của cô, rồi thở dài ngán ngẩm: ``Thế thì anh hỏi em cái phép tính
này:\( -3 - (-5)\). Tại sao em lại ra kết quả là -8?''

Cô em gái giật mình, ngẫm nghĩ lại một lúc, rồi trả lời anh trai với sự mập mờ: ``\dots\ Em
nghi là trừ một số âm thì giữ nguyên dấu trừ ạ?''

Nghe lời giải thích của cô, cậu vỗ thái dương, lắc đầu thất vọng. ``Kiến thức siêu cơ bản
đấy em ơi. Cái này anh dạy rất nhiều lần rồi mà còn chả nhớ.''

Chưa để Ryuuhou nói gì thêm, Kuon viết lại phép tính lên bảng, rồi từ tốn giảng giải lại
một lần nữa: ``Anh nhắc lại này: trừ của trừ là cộng. Thế này nhá: -3 trừ đi âm 5, nghĩa
là -3 cộng 5. Kết quả đúng là 2, không phải -8. Cái này anh dạy bao nhiều lần rồi đó, óc
em để trên mây à? Hay là não load (tải dữ liệu) chậm?''

Lời nói như mạt sát cô làm cho Ryuuhou nhăn nhó, bối rối đáp lại: ``\dots Em lại bị nhầm
thật rồi\dots'' Choco tiến tới, vội an ủi cô: ``Không sao đâu, Ryuuhou-sama. Lần sau em chỉ
cần cẩn thận hơn một chút thôi là sẽ ổn mà.''

A-chan nhìn cậu con trai đang hơi gắt gỏng với cô em gái, liền lên tiếng: ``Hikawa-san,
cậu nghĩ cậu đang có hơi khắt khe với em gái quá không?''

Kuon thở dài, đáp lại: ``Người nhà của em thì em càng phải nghiêm khắc chị ạ. Cũng sắp
sửa thi đại học đến nơi rồi đấy.''

Nhìn cô em gái đang thở dài chán chường vì bị anh trai nhắc nhở, Miko-chi ở ngay bên
cạnh cô, hớn hở chớp cơ hội: ``Haha! Hóa ra là Ryuuhou-chan cũng PON chẳng kém cạnh
gì Roboco đâu ha, ne!''

Nàng Người máy đang ỉu xìu ngồi ở bàn cuối, cảm thấy mình bị réo tên liền nói trong cơn
ngắc ngoải: ``Em không có PON thật mà\dots''

Ryuuhou quay phắt lại, mắt rực lửa phản ứng: ``Chị Miko-PON đừng có cười em! Chị
còn chẳng biết nhân chia trước cộng trừ sau cơ mà!''

Giờ thì đến lượt nàng Vu nữ nổi cáu khi bị cô ấy chê là PON: ``\textbf{Hả!? Cô em vừa nói cái gì cơ!? Có ngon thì nói lại xem nào!}''

Thấy hai người họ bắt đầu căng thẳng, Kuon gõ tay xuống bàn, cắt ngang: ``Thôi, thôi,
dừng. Hai người thôi đấu khẩu đi. Miko-chan, người ta nói đúng mà còn giãy nảy gì.''

*BỤP!* ``Hự!''

Choco-sensei cười nhẹ: ``Nếu xét về bối cảnh đó thì em thậm chí còn PON hơn cả
Ryuuhou-sama đấy. Sai kiến thức cơ bản đến thế mà.''

*BỤP!* ``Ặc!''

A-chan cũng nói góp: ``Có khác gì chó chê mèo lắm lông đâu. Không phải tự dưng người ta còn gọi cô là một trong những thành viên ngố nhất của công ty.''

*BỤP!* ``Ái\dots''

Miko-chi sau khi bị ba người họ buông ra những câu nói như tung đòn chí mạng vào niềm
kiêu hãnh của mình, cô như thể bị nốc-ao, nằm vật vờ xuống sàn, run lẩy bẩy: ``T-Tại sao
ba người lại có thể làm thế với tôi chứ, ne\dots''

Cậu quay sang cô em gái đang bật cười với hình dáng the thảm của nàng Vu nữ, rồi phán
xét với kết quả của cô: ``Về cơ bản, anh thấy em làm mấy câu tích phân, ma trận đều đúng
cả. Nhưng chỉ vì sai mấy câu cơ bản mà tụt hẳn 20 điểm đấy. Em được 62 điểm thôi.''

Cô em gái cười trừ, tỏ vẻ khá mãn nguyện dù miễn cưỡng phải chấp nhận kết quả này:
``\dots Ít nhất thì em vẫn sẽ được thưởng tiền.''

Mọi người cười phá lên trước phản ứng của cô em gái và nàng Vu nữ, biến phòng học trở
nên náo nhiệt hơn hẳn.

\ct

Như vậy, dưới đây là bảng tổng kết điểm của mười chín thí sinh đã tham dự cuộc thi này:

\begin{table}[H]
    \centering
    \begin{tabular}{|>{\raggedright\arraybackslash}p{0.25\linewidth}|>{\raggedright\arraybackslash}p{0.06\linewidth}|>{\raggedright\arraybackslash}p{0.06\linewidth}|>{\raggedright\arraybackslash}p{0.25\linewidth}|>{\raggedright\arraybackslash}p{0.06\linewidth}|>{\raggedright\arraybackslash}p{0.06\linewidth}|} \hline
        \textbf{Thành viên} & \textbf{Điểm} & \textbf{Loại} & \textbf{Thành viên} & \textbf{Điểm} & \textbf{Loại} \\ \hline
        Tokino Sora         & 86            & 1             & Nakiri Ayame        & 32            & 3             \\ \hline
        Roboco-san          & 0             & 4             & Oozora Subaru       & 30            & 3             \\ \hline
        Aki Rosenthal       & 74            & 2             & Ookami Mio          & 72            & 2             \\ \hline
        Akai Haato          & 82            & 1             & Sakura Miko         & 35**          & 3             \\ \hline
        Natsuiro Matsuri    & 54            & 2             & Nekomata Okayu      & 40*           & 3             \\ \hline
        Shirakami Fubuki    & 60            & 2             & Inugami Korone      & 60*           & 2             \\ \hline
        Yozora Mel          & 54            & 2             & Usada Pekora        & 64            & 2             \\ \hline
        Minato Aqua         & 58            & 2             & Shirogane Noel      & 28            & 4             \\ \hline
        Murasaki Shion      & 74            & 3             & Uruha Rushia        & 50            & 2             \\\hline
        Hikawa Ryuuhou      & 62            & 2             &                     &               &               \\ \hline
    \end{tabular}
    \caption{Chú thích: * Sau khi đã trừ 25\% điểm bài thi. ** Sau khi đã trừ 1 điểm.}
\end{table}

Korone nhìn bảng kết quả trên với vẻ mặt buồn thiu: ``Tức quá\dots\ Nếu không bị trừ chắc
tôi được nhận 10000 yên rồi\dots'' Okayu vỗ nhẹ vào vai cô bạn thân, nửa an ủi, nửa trêu
chọc: ``Thôi, còn đỡ hơn tôi, chuẩn bị chịu phạt này. May mà vẫn còn có chút tiền cầm
tay.''

Noel nhìn cặp đôi rồi quay sang Ryuuhou, ánh mắt có phần ganh tị: ``Không công bằng!
Em gái của Kuon-senpai được nhận tiền kìa!''

Cô em gái của cậu Tổng khẽ nhún vai, giọng bình thản nhưng không giấu được chút tự
hào: ``Em cũng chỉ nhận được có năm nghìn thôi, chị đừng có cảm xúc thế, Noel-onee-chan. Với lại, chị
bị 28 điểm thì trách ai?''

Roboco ngồi phịch xuống ghế, tay chống cằm, giọng uể oải: ``Bọn chị còn chẳng được xu
nào nữa này! Tại sao chị lại có thể làm sai được chứ\dots''

Miko nhìn nàng Người máy đang ỉu xìu đó, bật cười: ``Bà bất cẩn thì bà chết thôi, ne. Tôi
mà không phải ưu tú thì chắc cũng thế!''

Trong khi đó, Shion liếc nhìn xấp tiền trong tay cậu Tổng, ánh mắt toát ra vẻ lấp lánh:
``Vậy anh đã chuẩn bị chưa? Trông xấp tiền kia nặng tay phết nhỉ!''

Kuon mỉm cười, xòe tiền từ trong tay ra: ``Luôn sẵn sàng. Nào, mọi người lại đây nhận
phần thưởng đi.''

Cậu bắt đầu chia tiền cho mọi người: Sora và Haato, mỗi người nhận được tờ mười nghìn
yên, khuôn mặt cả hai ánh lên niềm vui rạng rỡ.

Tiếp theo, Aki, Fubuki, Matsuri, Mel, Shion, Mio, Ryuuhou, Korone và Pekora mỗi người
nhận năm nghìn yên. Korone dù buồn nhưng vẫn cảm ơn rối rít. Rushia cũng nhận tờ
năm nghìn yên, khẽ cúi đầu, mỉm cười với ánh mắt long lanh.

Riêng với cô em gái của cậu, Kuon không quên nhắc: ``Tiêu tiết kiệm vào nha. Tiền đưa cho em toàn thấy tiêu hết sạch!'' Tất cả mọi người khẽ cười với lời nhắc nhở với người nhà của
cậu.

Kuon giơ tay lên, ra hiệu cả lớp chú ý: ``Tất nhiên, những người thuộc loại 3 cũng được
nhận giải an ủi nữa.''

Aqua đang ngồi thất thần, chuẩn bị chịu phạt từ nàng Thần tượng và nàng Song tính,
ngẩng đầu lên với vẻ ngạc nhiên: ``Bọn em cũng được nhận tiền ạ!?''

Cậu gật đầu, xòe thêm vài tờ tiền mệnh giá thấp hơn nữa: ``Ừ. Bình đẳng mà. Dù sao
cũng có cố gắng làm bài. Mỗi người 2000 yên.'' Cậu phát tiền cho Miko, Aqua, Ayame,
Subaru và Okayu. Dù là số tiền không lớn, nhưng tất cả đều cảm ơn vì ít nhất cậu còn nghĩ tới họ.

Nàng Miêu thậm chí còn giơ tiền lên, cười híp mắt: ``Cảm ơn anh nha\~{}'' rồi quay sang
chung vui với Korone.

Nhìn đa số mọi người được lì xì, Roboco với vẻ mặt như muốn đòi quyền lợi, than thở:
``Còn bọn em thì sao? Chẳng lẽ được số 0 là không có gì luôn ạ!?''

Noel gật đầu đồng tình, vẻ mặt vẫn còn uất ức: ``Ít nhất anh cũng phải có chút gì an ủi
bọn em chứ, Kuon-senpai!''

Pekora nhìn hai người họ đòi hỏi cậu như vậy, mỉa mai: ``Hai người không làm mà đòi có
ăn cơ à? Thế thì chỉ có đi cạp đất mà ăn thôi, peko!''

Kuon nhìn phản ứng bất mãn của hai cô nàng loại bốn, bật cười: ``Hai người thì, ờm\dots\ xin Sora-chan với Haachama đi nhé. Phần thưởng của anh coi như xong hết rồi, còn lại
chỉ còn phạt thôi.''

A-chan đứng bên cạnh, bật cười và bổ sung thêm: ``Tôi sẽ ghi nhận đề nghị này, nhưng
để lần sau nhé. Còn giờ thì chuẩn bị chịu phạt đi!''

Hai người lầm bầm, nhưng nhìn vẻ mặt nghiêm nghị của cậu con trai và A-chan, họ đành
phải im lặng chịu trận.

``Được rồi, Sora-chan, Haato-chan, em muốn những người loại ba làm gì?'' cậu hỏi, ánh
mắt chờ đợi một câu trả lời không quá\dots\ nguy hiểm hay đáng sợ.

Đùng một cái, mặt của Sora tối sầm lại, đôi mắt bỗng chốc mất đi vẻ long lanh thường
thấy, thay vào đó là ánh nhìn ``nguy hiểm'' đầy ẩn ý, như ánh mắt của một chú mèo đang
vờn mồi: ``Họ có thể mua vui cho chúng ta được mà\dots\ nhỉ?''

Mọi người trong phòng rùng mình với phản ứng đáng sợ của nàng Thần tượng.
``Đó\dots\ cũng là một cách - Mà Sora-chan, trông mặt em nguy hiểm quá\dots'' cậu con trai
toát mồ hôi hột.

A-chan thì cũng sởn gai ốc không kém, vội vàng nói với người bạn của mình: ``Bình tĩnh
lại đi nào, Sora\dots''

Tất cả mọi người nhìn cô nàng đang tỏa ra sát khí kia, nuốt nước bọt trong lo sợ. Đến cả
những thành viên thuộc loại hai - những người đã thoát được hình phạt - cũng khẽ run rẩy
sợ hãi.

\begin{dialogue}
    \speak{Rushia} S-Sora-senpai đáng sợ quá-nanodesu\dots
    \speak{Subaru} Chị ấy có vẻ mặt khó tả thật\dots
    \speak{Miko} B-Bà đang định làm gì chúng tôi, ne!?
\end{dialogue}

Đột ngột, cô nàng ấy quay ngoắt 180 độ, nụ cười sáng bừng như ánh nắng ban mai: ``Nói
thế thôi, em nghĩ họ có thể làm những hoạt động mà ngày Tết ở đền chùa bọn anh hay
làm ấy.'' Tất cả mọi người đều thở phào nhẹ nhõm với yêu cầu không quá điên rồ này.

``À, kiểu như giúp việc ở đền chùa? Phụ dọn dẹp, chuẩn bị đồ lễ hay hỗ trợ khách thập
phương ấy à?'' Kuon hỏi xác nhận lại.

``Chính xác!''

Cậu con trai gật gù, quay sang nhìn cô nàng Song tính đang suy nghĩ. ``Haato-chan thì
sao? Có ý tưởng gì khác không?''

Haato đáp, dường như không có ý tưởng gì khác: ``Em cũng đồng quan điểm với chị ấy.
Làm mấy việc đơn giản thôi, như phát bùa chúc phúc hay giúp bày đồ cúng chẳng hạn.''

Cậu cười, tay gõ nhẹ lên tấm bảng kế hoạch: ``Vậy tức là tí nữa họ sẽ\dots\ kiểu, làm người
hỗ trợ ở đền chùa trong vai trò thực tập sinh à? OK. Ý tưởng tốt đấy.''

Phản ứng với hình phạt mà Sora và Haato của những người thuộc loại ba cũng khá khác
nhau. Trong khi Okayu, Ayame và Subaru và Miko tỏ vẻ khá thích thú với công việc này,
chỉ có Pekora thì nghi ngờ về ý tưởng này.

\begin{dialogue}
    \speak{Pekora} Đền chùa thì có liên quan gì đến việc kiểm tra đâu, peko?
    \speak{Subaru} Thôi đi, ít ra thì không phải nhảy múa hay làm trò hề là tốt rồi!
    \speak{Aqua} Em còn tưởng là chị ấy bắt chúng ta phải làm trò gì quái gở hơn cơ\dots
\end{dialogue}

Tiếng cười bật lên, nhưng mọi người đều hiểu, một ``nhiệm vụ ngày Tết'' đang chờ đợi
những thành viên ``may mắn'' ở loại ba!

``Giờ đến hình phạt của hai đứa đại PON kia.''

Câu nói ấy vừa vang lên, tất cả ánh mắt lập tức đổ dồn về phía Noel và Roboco, hai người
đang ngồi rìa góc phòng, cố gắng trốn tránh sự chú ý.

Cảm nhận được khí lạnh đột ngột tràn qua, họ từ từ ngẩng đầu lên, và đập vào mắt họ là
hàng loạt gương mặt đầy ``sát khí'' của các thành viên khác đang tiến tới.

``Ê-tồ\dots\ Mọi người đang làm gì vậy?'' Roboco e sợ hỏi, cố rướn người ra phía sau. Shion
- một trong hai người nghĩ ra hình phạt bá đạo này - mỉm cười đáp: ``Phạt chứ làm gì.''

Noel - người còn lại nghĩ ra ý tưởng phạt đó - định cất lời phản đối nhưng chưa kịp mở
miệng đã bị Pekora nhanh chóng lao tới: ``Có nhất thiết phải sồn sồn thế không- Á!!''

\ct

Chỉ trong nháy mắt, cả hai cô nàg bị các cô gái còn lại kéo ra giữa phòng. Không kịp phản
kháng, họ bị ép phải đeo lên cổ hai tấm bảng lớn ghi dòng chữ đỏ chói: ``\textbf{Tôi dốt toán}''
bằng cả tiếng Nhật và tiếng Etsunan.

Nàng Người máy cố gắng phản kháng: ``K-Khoan đã! Sao mọi thứ lại thành ra thế này!?''
Noel cũng cố gắng phản kháng không kém: ``Em cũng đâu đến mức ấy đâu chứ! Thả
mình ra!!''

Mặc kệ tiếng kêu gào, nhóm ``ban phạt'' dẫn đầu bởi Shion và Rushia đã quyết định nâng
hình phạt lên một tầm cao mới.

Họ xếp thành đội hình ``diễu binh,'' kéo hai nạn nhân đi dọc hành lang của trường đại học
mà Kuon đang học. Theo sau hai ``phạm nhân'' là các thành viên còn lại của \hll\ và
cô em gái Ryuuhou, theo sau cùng là A-chan và Kuon.

Nàng Phù thủy đứng đầu, reo lên đầy phấn khích: ``Mọi người ơi, xem này! Có hai tên
dốt toán đây, bà con ơi!''

Bên cạnh cô, nàng Pháp sư cũng hớn hở không kém: ``Chúng tôi có hai kẻ dốt toán nhất
trong nhóm này-nanodesu!''

Hai nạn nhân khốn khổ, một người cố van xin, người còn lại thì giãy nảy:

\begin{dialogue}
    \speak{Roboco} *rơi nước mắt anime* Tha cho bọn chị đi mà!
    \speak{Noel} Tại sao mọi chuyện lại xảy ra với mình chứ\dots?
\end{dialogue}

Pekora ngay sau hai người họ, mỉm cười châm chọc: ``Bà là người đề xuất ý tưởng hình
phạt đấy chứ còn gì, peko!'' Fubuki nhoẻn miệng, nói với tông giọng mỉa mai: ``Thật
luôn. Cái này gọi là gậy ông đập lưng ông đấy!''

Cả Noel và Roboco lúc này cũng chỉ có thể ngước đầu lên thất thanh đến vọng cả trời:
``\textbf{K-Không thể nào!!}''

Rất may mắn cho họ, vào ngày mùng Ba, trường vắng tanh vì phần lớn sinh viên còn đang
nghỉ Tết. Ngoài một vài người có đến chúc Tết thầy cô giáo, chỉ còn lại nhóm bạn cùng
nhau vui đùa tại đây.

Nhưng điều trớ trêu là, nhóm sinh viên ngành khoa học máy tính vô tình bắt gặp ``cuộc
diễu binh'' của các cô gái và lập tức bật cười nghiêng ngả, không quên quay chụp lại cảnh
tượng hiếm có này.

Bản thân Rushia reo lên thích thú: ``Được lắm, để tôi đưa lên Twitter nhé. Xem ai dám
quên ngày này!'' Pekora cười khì: ``Khỏi cần đâu, đằng nào A-chan cũng sẽ quay lại
khung cảnh này mà, peko!''

Cả nhóm bật cười, trong khi hai nhân vật chính chỉ biết cúi đầu chịu trận.

Kuon và A-chan đứng quan sát cả nhóm đang pha trò sau bài kiểm tra toán. Tiếng cười
vang khắp căn phòng, trong khi hai người quản lý chỉ có thể nhìn khung cảnh hỗn loạn
đó bằng những cái cười trừ.

``Cũng may họ là các idol đấy,'' A-chan mở lời, quay sang cậu Tổng. Kuon đáp lại: ``Vâng.
Nhưng mà chẳng biết họ còn có thể giữ được hình tượng ấy đến bao giờ nữa\dots''

Họ tiếp tục xem mọi người đang trêu đùa với nhau, không nói một lời nào. Sau một thoáng
im lặng, nàng Đeo kính tiếp tục quay sang cậu: ``Cậu sống ở đây chắc áp lực lắm nhỉ?''

Cậu đáp, đôi mắt tỏ vẻ đượm buồn: ``Ở đất Bắc này thì sống chẳng dễ chút nào, toàn học
với học thôi ạ. Người ta coi trọng tri thức cực kỳ, nên ai cũng phải gồng mình cố hết sức thôi ạ.''

``Nghe cũng áp lực như ở Tokyo vậy. Nhìn cậu thế này, tôi hiểu cậu đã phải phấn đấu thế
nào để được như bây giờ.''

Kuon mỉm cười gật đầu nhẹ: ``Có lẽ còn hơn cả chị nghĩ nữa đấy ạ. Mà, em cảm ơn chị
vì đã quan tâm tới em.''

Cậu bất chợt nhớ ra về điều mà cậu đã từng nghe trước đây: ``Mà, nhắc đến đất Bắc thì,
em nhớ Ayame-chan từng nói em ấy đã ăn Tết một lần ở Etsunan, phải không?''

``Ừ, cô ấy có kể. Ở trong phía Nam.''

``Vậy chắc là Ayame-chan đã hơi sốc văn hóa hoặc cảm thấy quen thuộc đến mức chẳng
buồn nói. Hai vùng ấy khác nhau hoàn toàn luôn, chị ạ.''

A-chan bật cười nhẹ, gật gầu đồng tình: ``Đến mức người miền ngoài ra miền trong, hoặc
ngược lại, có cảm giác như đi nước ngoài ấy nhỉ? Tôi cũng được nghe nói rồi.''

Nửa ngày trôi qua một cách nhẹ nhàng. Buổi trưa, cả nhóm về nhà quây quần bên đống
bánh kẹo và mứt Tết - những món ăn vặt họ đã thưởng thức suốt ba ngày qua. Không ai
ăn cơm, tất cả đều chuẩn bị năng lượng cho các hoạt động thú vị vào buổi chiều.

\end{document}